\documentclass[a4paper, 12pt]{report}

%package
\usepackage[french]{babel}
\usepackage[margin=0.8in]{geometry}
\usepackage{graphicx} % Required for inserting images
\usepackage{afterpage}
\usepackage{amssymb}
\usepackage{amsmath}
\usepackage{amsthm}
\usepackage{mathabx}
\usepackage{tikz}
\usepackage{thmtools}
\usepackage{framed}
\usepackage[hidelinks]{hyperref}

%% préambule
\usepackage[
   backend=biber,        % compilateur par défaut pour biblatex
   sorting=nyt,          % trier par nom, année, titre
   citestyle=authoryear, % style de citation auteur-année
   bibstyle=alphabetic,  % style de bibliographie alphabétique
]{biblatex}
\addbibresource{ref.bib}
%% corps du document

\newcommand{\R}{\mathbb{R}} % : pour l'ensemble de réels
\newcommand{\N}{\mathbb{N}} % : pour l'ensemble des entiers naturels
\newcommand{\Z}{\mathbb{Z}} % : pour l'ensemble des entiers relatifs
\newcommand{\Q}{\mathbb{Q}} % : pour l'ensemble des rationnels
\newcommand{\C}{\mathbb{C}} % : pour l'ensemble des complexes
\newcommand{\F}{\mathbb{F}} % : pour les corps finis

\newcommand{\iif}{\Leftrightarrow} % : pour le si et seulement si (symbole équivalence)
\newcommand{\ra}[1][]{\overset{#1}{\rightarrow}} % : pour la flèche vers la droite
\newcommand{\Ra}{\Rightarrow} % : pour le symbole implication
\newcommand{\Lim}[2]{\underset{#1 \ra #2}{\lim}}
\renewcommand{\.}{\cdot} % : pour le point

\newcommand\emptypage{
    \null
    \thispagestyle{empty}
    
    \newpage
    }

\usepackage{thmtools}
    \declaretheoremstyle[
    headfont=\bfseries, 
    notebraces={}{},
    headformat=Preuve : \NOTE,
    bodyfont=\normalfont\itshape,
    headpunct={},
    spacebelow=\parsep,
    spaceabove=\parsep,
    qed=\qedsymbol
]{proofstyle}
\declaretheorem[name=Démonstration, numbered=no, style=proofstyle]{preuve}
\declaretheoremstyle[
    headfont=\bfseries, 
    notebraces={}{},
    headformat=Théorème: \NOTE,
    bodyfont=\normalfont\itshape,
    headpunct={},
    postheadspace=\newline,
    %postheadhook={\textcolor{black!20}{\rule[.3ex]{\linewidth}{0.4pt}}\\},
    spacebelow=\parsep,
    spaceabove=\parsep,
    mdframed={
        backgroundcolor=gray!5, 
            linecolor=black, 
            linewidth=1pt,
            innertopmargin=6pt,
            roundcorner=5pt, 
            innerbottommargin=6pt, 
            skipbelow=\parsep, 
            skipabove=\parsep } 
]{theoremstyle}
    
\declaretheorem[name=Théorème, numbered=no, style=theoremstyle]{theorem}

\newtheoremstyle{break}%
{}{}%
{\itshape}{}%
{\bfseries}{}% % Note that final punctuation is omitted.
{\newline}{}



\theoremstyle{break}
\newtheorem{definition}{Définition}
\newtheorem{proposition}{Proposition}

\theoremstyle{remark}
\newtheorem{remark}{Remarque}
\newtheorem{exemple}{Exemple}

\theoremstyle{proofstyle}

\title{La transformation de Weyl}
\author{Noirot Lucas \\ Encadrant : Giuseppe Dito}

\usepackage{fancyhdr}
\pagestyle{fancy}

\AtEndDocument{\label{lastpage}}

\fancyhead[L]{Noirot Lucas}\fancyhead[C]{}\fancyhead[R]{Université de Bourgogne}
\fancyfoot[C]{\begin{tabular}{cc} La transformation de Weyl \\ \textbf{Page \thepage/\pageref{lastpage}} \\ \end{tabular}}

\renewcommand{\footrulewidth}{0.4pt}

\date{07 mai 2023}

\setlength{\parindent}{0pt}
\begin{document}


\maketitle

\emptypage

\chapter*{\centering Introduction}


En \textbf{mécanique classique}, l'état d'un système constitué d'une particule \textbf{se mouvant sur une droite} à un instant $t$ peut-être décrit par la donnée de la \textbf{position} $x$ et de la \textbf{vitesse} $v$. On peut prédire son évolution en appliquant l'équation de Newton $\vec{a} = \frac{\vec{F}}{m}$, où $\vec{F}$ représente la somme des forces extérieures agissant sur ce système, $m$ est la masse de la particule et $\vec{a}$ son accélération. L'étude de l'évolution d'un objet dans $\R^n$ se fait souvent dans l'espace des phases $\R^n \times \R^n = \{(\xi,x)$ tel que $ \xi \in \R^n, x \in \R^n\}$, l'espace dans lequel tous les états possibles d'un objet sont représentés, où en général $x$ représente la position et $\xi$ l'impulsion, les lettres $q$ et $p$ seront aussi utilisées. Par exemple on peut choisir de représenter une balle de tennis, assimilée à un point matériel, par sa position et sa vitesse dans l'espace, donc par un point de $\R^3 \times \R^3$ à chaque instant $t$ ! Puis en appliquant l'équation de Newton on va pouvoir prédire l'évolution au cours du temps de la position de la balle de tennis. On appelle grandeur physique toute propriété d'un système qui peut être mesurée ou calculée et dont les valeurs possibles sont réelles. Ce sont des fonctions définies sur l'espace des phases, par exemple l'énergie du système considéré.\vspace{7pt}

En \textbf{mécanique quantique}, c'est conceptuellement différent, on ne décrit plus l'état d'un système par sa position et son impulsion mais par \textbf{une fonction à valeurs complexes appelée fonction d'onde} $\Psi_t$ à chaque instant $t$.  On parle d'état quantique. Les fonctions d'onde sont régies par la célèbre équation de Schrödinger $i \hbar \frac{d \Psi}{dt}=\hat{H}\Psi$, où $\hat{H}$ s'appelle l'Hamiltonien du système et joue le même rôle que $\vec{F}$, c'est un opérateur qui agit sur les fonctions d'onde $\psi _t$, $\hbar$ vaut $\frac{h}{2 \pi}$, $h$ étant la célèbre constante de Planck. On veut maintenant manipuler ces fonctions d'onde à l'aide d'opérateurs auto-adjoints qui agissent sur ces états. Ces opérateurs sont appelés des \textbf{observables} et jouent un rôle analogue à celui des grandeurs physiques en mécanique classique. Ils sont appelés observables car c'est par l'action de ces opérateurs sur les états qu'on peut comprendre et observer les particules quantiques. Par exemple, une observable peut être l'opérateur représentant la position, associe la valeur moyenne de la position dans un état $\psi$.\vspace{7 pt}

Deux opérateurs sont \textbf{au centre de la mécanique quantique}, l'opérateur de \textbf{position} $\hat{Q} : \psi \mapsto x \psi$ défini sur $\mathcal{H}=L^2(\R^n)$ et l'opérateur d'\textbf{impulsion} $\hat{P} : \psi \mapsto -i \hbar \frac{\partial \psi}{\partial x}$ défini sur $\mathcal{H}=L^2(\R^n)$. Ces espaces seront amenés à changer lors de ce mémoire.\vspace{7 pt}


Lorsqu'on mesure une observable sur un état quantique, on obtient une des valeurs spectrales de l'observable avec une certaine probabilité. Cette probabilité est donnée par le carré du module de la projection de l'état quantique sur le vecteur propre associé à cette valeur propre. Après la mesure, l'état quantique du système est projeté sur ce vecteur propre. C'est ce qu'on appelle le postulat de la réduction du paquet d'ondes. Il est important de noter que \textbf{les observables ne commutent pas toujours entre elles}, ce qui a pour conséquence le fameux principe d'incertitude de Heisenberg, donné par la célèbre inégalité $\Delta_x \Delta_p \geq \frac{\hbar}{2}$, où $\Delta_x$ (resp. $\Delta_p)$ représente l'écart-type de $\hat{Q}$ (resp. $\hat{P}$). Ce principe formulé en 1927 par Werner Heisenberg, stipule qu'il est impossible de mesurer simultanément avec une précision arbitraire la position et l'impulsion d'une particule quantique. En d'autres termes, plus on mesure précisément la position d'une particule, moins on peut connaître précisément son impulsion, et inversement. Ce principe se généralise pour tout couple d'observables ne commutant pas. Ce qui rend \textbf{difficile} la formulation de cette théorie dans l'espace des phases. Ce principe est souvent interprété comme une limitation fondamentale de la nature de la réalité physique quantique. \vspace{7 pt}

C'est dans ce contexte que la transformée de Weyl trouve toute son \textbf{importance}. \vspace{7 pt}

Dans ce mémoire, on sera intéressé par l'étude d'une particule se déplaçant dans $\mathbb{R}^n$, on prendra l'espace des opérateurs de $L^2(\R^n)$ auto-adjoints. L'une des motivations de la transformée de Weyl est de vouloir formuler la mécanique quantique sur l'espace des phases, qui de base est formulée sur l'espace des positions, avec la grandeur physique $q$ (resp. $p$). \vspace{7 pt}

La \textbf{problématique de la quantification} est d'associer à chaque fonction $f$ sur l'espace des phase $\R^{2n}$ un opérateur $\hat{f}$ sur $L^2(\mathbb{R}^n)$ \textbf{auto-adjoint}. La première procédure de quantification et l'une des plus satisfaisante a été proposée par Weyl, elle propose d'associer un opérateur $\sigma(\hat{P},\hat{Q})$ agissant sur $\mathcal{H} = L^2(\R^n)$ à une fonction $\sigma(\xi, x)$ par la formule

$$\sigma(\hat{P},\hat{Q}) = \frac{1}{2 \pi \hbar} \iint \hat{\sigma} (p,q) e^{ \frac{i}{\hbar}(p\hat{P}+q\hat{Q})} dpdq$$

où $\hat{P}, \hat{Q}$ représentent les opérateurs d'impulsion et de position, et $\hat{\sigma}$ est la transformée de Fourier $\sigma$. On appellera cette formule $(Weyl)$ dans la suite du mémoire.

Cette formule vient du fait suivant :
Si on se donne une fonction $\sigma(\xi, x)$, on aimerait lui faire correspondre un opérateur $\sigma(D,x)$ auto-adjoint. En supposant que l'on peut écrire $\sigma(x, \xi)$ en son écriture de Fourier $\iint \hat{\sigma} (p,q) e^{2 \pi i (p \xi + q x)} dpdq$, en supposant que $e^{2 \pi i (p \xi + qx)}, p, q \in \R^n$ devrait correspondre à l'opérateur $e^{2 \pi i (p\hat{P} + q\hat{Q})}$. On serait alors amené à postuler la définition $(Weyl)$. \vspace{10 pt}

Le but de ce mémoire est d'expliquer la formule $(Weyl)$ mathématiquement et physiquement, qu'on illustrera par quelques exemples. On finira par étudier ses propriétés, en faisant un lien solide entre les aspects physique et mathématique des objets que l'on va être amenés à étudier au cours de ce mémoire afin d'expliquer $(4.2)$. On a besoin de cette correspondance car il n'est pas si évident d'associer un opérateur à une fonction sur l'espace des phase. Il serait tentant de penser que $pq$ devrait être associé à $\hat{P}\hat{Q}$ mais $pq = qp$ alors que le commutateur de $\hat{P}$ et $\hat{Q}$, $[\hat{P}, \hat{Q}] \neq 0$, c'est non commutatif. On doit alors réfléchir à une manière plus subtile de faire cette correspondance. On devra alors finalement munir cet espace de phase de nouvelles opérations afin de garder une belle structure.\vspace{10 pt}



La première partie sera consacrée à l'étude de certains objets mathématiques nécessaires à la compréhension de la suite du mémoire, elle introduit les bases mathématiques de la mécanique quantique ainsi que ses postulats, puis on va étudier la transformée de Fourier et son interprétation physique qui sera au c\oe ur de beaucoup d'objets étudiés dans ce mémoire, on finira ensuite cette première partie par une introduction aux opérateurs pseudo-différentiels ainsi que leurs symboles, objet mathématique au c\oe ur de la transformée de Weyl, bien que non nécessaire mais fournissant de mon point de vue une bonne motivation à l'étude de la transformée de Weyl. \vspace{10pt}

La deuxième partie, centre de ce mémoire, sera consacrée à l'étude de la transformée de Weyl, on y commencera par étudier la transformée de Fourier-Wigner et son lien avec le groupe de Heisenberg et ses représentations irréductibles unitaires, puis on introduira la transformée de Wigner, qui sera avant tout expliquée pour motiver la transformée de Weyl rigoureusement, on finira cette partie par l'illustration de quelques exemples de correspondance entre fonction et opérateur par cet outil qu'est la transformée de Weyl, puis par le développement d'un calcul fonctionnel et agréable sur ces nouveaux objets. \vspace{10pt}
\emptypage

\tableofcontents


\emptypage
\emptypage

\chapter{Mathématiques de la mécanique quantique}

\section{Rappels sur les espaces de Hilbert et les opérateurs}
Cette partie servira de rappels mathématiques sur des notions très utiles dans ce mémoire.

On introduit la notion d'espace de Hilbert, un objet qui est \textbf{à la la base de la formulation} de la mécanique quantique.

\begin{definition}[Espace de Hilbert]
    On dit que $\mathcal{H}$ est un \textbf{espace de Hilbert} si c'est un espace vectoriel complexe, muni d'un produit scalaire hermitien, noté $\langle \cdot, \cdot \rangle$ : $\mathcal{H} \times \mathcal{H} \rightarrow \mathbb{C}$, on le considère anti-linéaire dans le premier argument. Et si il est complet pour la norme associée à son produit scalaire.
\end{definition}


\begin{remark}
    Un espace de Hilbert n'est pas nécessairement de dimension finie.
\end{remark}

\begin{definition}[Base hilbertienne d'un espace de Hilbert]
    Soit $\mathcal{H}$ un espace de Hilbert, on dit que $\{e_n\}_{n \in \mathbb{N}}$ est une \textbf{base hilbertienne} de $\mathcal{H}$ si :
    $\langle e_i, e_j \rangle = \delta_{ij}$ et pour tout $n \in \mathbb{N}$, $\langle \psi, e_n \rangle = 0$ entraîne $\psi = 0$.
\end{definition}

Deux inégalités seront maintenant présentée, celle de Cauchy-Schwarz pour les espaces $L^p$, et celle de Young pour le produit de convolution. Elles seront très utiles pour les démonstrations du mémoire.

\begin{proposition}[Inégalité de Cauchy-Schwarz]
    Soit $\mathcal{H}$ un espace de Hilbert muni du produit scalaire $\langle \., \. \rangle$. Alors pour tout $x$, $y \in \mathcal{H}$,
    $$| \langle x, y \rangle | \leq \| x \| \| y \|.$$
\end{proposition}

\begin{preuve}
    Soit $x, y \in \mathcal{H}$.

    Si $\| x \| = 0$ alors $x=0$ et l'inégalité est immédiate.

    Si $\| x\| \neq 0$ alors en développant
    $$\|y-\langle x, y \rangle \frac{x}{\|x\|^2}\|^2 \geq 0$$
    Il vient
    $$\|x\|^2\|y\|^2 \geq \langle x, y \rangle ^2$$
    Comme $\|x\| \|y\| \geq 0$, on a finalement $\|x\| \|y\| \geq |\langle x, y\rangle|$
\end{preuve}

Dans la proposition suivante, les espaces $L^p$ désigneront les espaces de Lebesgue.

\begin{proposition}[Inégalité de Young pour le produit de convolution]
    Soit $f$ dans $L^p$ et $g$ dans $L^q$, soit $p, ~q,~r \in [1, \infty ]$ tel que $\frac{1}{p} + \frac{1}{q} = 1 + \frac{1}{r}.$

    Alors le produit de convolution $f * g$ appartient à $L^r$ et,
    $$\|f * g\|_r \leq \|f\|_p\|g\|_q.$$
\end{proposition}

Cette inégalité sera assumée pour vraie, la preuve utilisant deux autres inégalités qu'il aurait fallut introduire, ce qui n'est pas le sujet de ce mémoire. \vspace{5pt}

L'autre objet \textbf{fondamental} de la mécanique quantique est la notion d'opérateur.
\begin{definition}[Opérateur]
    Un \textbf{opérateur} $T$ entre deux espaces de Hilbert est une application $T : \mathcal{D}(T) \subset \mathcal{H}_1 \ra \mathcal{H}_2$ telle que,
    \begin{itemize}
        \item $\mathcal{D}(T)$ est un sous-espace vectoriel de $\mathcal{H}_1$.
        \item T est une application linéaire : $T(\alpha u + \beta v) = \alpha T u + \beta T v$ pour tout $\alpha$, $\beta \in \mathbb{C}$, pour tout $u$, $v \in \mathcal{D}(T)$.
    \end{itemize}
    où $\mathcal{D}(T)$ est le domaine de T.

    Quand $\mathcal{H}_1 = \mathcal{H}_2 = \mathcal{H}$ on dit que $T$ est un \textbf{opérateur sur $\mathcal{H}$}.
\end{definition}

On considère habituellement les opérateurs avec un \textbf{domaine dense}.

\begin{definition}[Opérateur borné]
    Un opérateur $T : \mathcal{D} (T) \subset \mathcal{H}_1 \ra \mathcal{H}_2$ est \textbf{borné} si il existe une constante $c > 0$ telle que
    $$\|T u\|_2 \leq c \|u\|_1, \text{ pour tout }u \in \mathcal{D}(T).$$

    La norme d'un opérateur borné est la plus petite constante c vérifiant l'inégalité précédente, et elle est donnée par
    $$\| T \| = \underset{u \in \mathcal{D} (T), u \neq 0}{sup} \frac{\|T u\|_2}{\|u\|_1}.$$
\end{definition}

\begin{proposition}
    Un opérateur est borné si et seulement si il est continu.
\end{proposition}\vspace{5pt}

\begin{theorem}[Théorème de Riesz]
    Soit $l : \mathcal{D}(l) \subset \mathcal{H} \ra \mathbb{C}$ une forme linéaire densément définie. Alors il existe un unique vecteur $v_l \in \mathcal{H}$ tel que,
    $$l(u) = \langle v_l, u \rangle, \text{ pour tout } u \in \mathcal{D}(l).$$

    On a $\|l\|=\|v_l\|$.
\end{theorem}\vspace{5pt}

On sera souvent amenés dans ce mémoire à prolonger des opérateurs $\hat{A}$ densément définis sur des espaces plus grand, souvent pour généraliser le résultat aux espaces $L^2(\R^n)$, on doit alors savoir qu'il y a unicité du prolongement. \vspace{7pt}


On introduit maintenant la notion d'\textbf{adjoint d'un opérateur}, notre but est d'associer à un opérateur $T$ densément défini un opérateur $T^*$ de domaine $\mathcal{D}(T^*)$ qui a la propriété $\langle T^*, u \rangle = \langle v, T u \rangle$ pour tout $u \in \mathcal{D} (T)$, pour tout $u \in \mathcal{D}(T^*)$. \vspace{5pt}

On définit $\mathcal{D}(T^*) := \{ v \in \mathcal{H} ; u \mapsto \langle v, T u \rangle$ est une forme linéaire bornée sur $\mathcal{D}(T)\}$, d'après le théorème de Riesz, il existe un unique $v^* \in \mathcal{H}$ telle que $\langle v^*, u \rangle = \langle v, Tu \rangle$.

\begin{definition}[L'adjoint d'un opérateur]
    Soit $T : \mathcal{D}(T) \subset \mathcal{H} \ra \mathcal{H}$ un opérateur densément défini. L'\textbf{adjoint} de $T$ est l'opérateur $T^* : \mathcal{D} (T^*) \subset \mathcal{H} \ra \mathcal{H}$ définit de façon \textbf{unique} par les conditions :
    \begin{itemize}
        \item $\mathcal{D}(T^*) := \{ v \in \mathcal{H} ; \exists ! v^* \in \mathcal{H}$ tel que $\langle v^*, u \rangle = \langle v, T u \rangle \text{ pour tout } u \in \mathcal{D}(T) \}$.
        \item $T^*v= v^*$ pour tout $v \in \mathcal{D}(T^*)$.
    \end{itemize}
\end{definition}

\begin{definition}[Opérateur auto-adjoint]
    Un opérateur densément défini $T : \mathcal{D}(T) \subset \mathcal{H} \ra \mathcal{H}$ est \textbf{auto-adjoint} si $T^* = T$.
\end{definition}
En particulier, on a $\mathcal{D} (T) = \mathcal{D}(T^*)$.

\begin{definition}[Opérateur symétrique]
    Un opérateur $T : \mathcal{D} \subset \mathcal{H} \ra \mathcal{H}$ est \textbf{ symétrique } si
    $\langle u, Tv \rangle = \langle T u, v \rangle, ~~\text{ pour tout } u, ~v \in \mathcal{D}(T).$
\end{definition}

En conséquence, pour les opérateurs auto-adjoints, on a,
$$\langle u, T v \rangle = \langle T u, v \rangle, ~~\text{ pour tout } u,~v \in \mathcal{D}(T).$$

En particulier, un opérateur auto-adjoint est symétrique.\vspace{10 pt}

%\begin{definition}[Spectre et valeurs propres]
%    Soit $T : \mathcal{D} \subset \mathcal{H} \ra \mathcal{H}$
%    \begin{itemize}
%        \item Un nombre complexe $\lambda$ est une valeur propre de $T$ s'il existe un vecteur $h \in \mathcal{H}, h \neq 0$ tel que $Th = \lambda h$. Un tel %vecteur est appelé un \textbf{vecteur propre} de $T$ associé à la \textbf{valeur propre} $\lambda$.
%        \item On appelle \textbf{spectre ponctuel} de $T$, et on note $\sigma _p(T)$ l'ensemble des valeurs propres de $T$.
%        \item Le \textbf{spectre} de $T$ est défini comme
%            $$\sigma (T) := \{ \lambda \in \mathbb{C} \text{ tel que } T _ \lambda Id \text{ n'est pas inversible.} \}$$
%    \end{itemize}
%\end{definition}

En général en dimension infinie, le spectre d'un opérateur \textbf{n'est pas égal} à l'ensemble de ses valeurs propres, on a une inclusion de l'ensemble des valeurs propres dans le spectre. Parce qu'un opérateur inversible est nécessairement injectif.
L'égalité est vraie en dimension finie grâce au théorème spectral et au théorème du rang. Il faudra donc faire attention lorsqu'on parle du spectre ou des valeurs propres d'un opérateur.\vspace{7pt}

Une propriété intéressante des opérateurs auto-adjoint est par rapport à leur spectre.
\begin{proposition}
    Le spectre d'un opérateur auto-adjoint est \textbf{réel}.
\end{proposition}

Toutes ces propositions sur les espaces de Hilbert ou les opérateurs ont été assumées pour vrai, l'objectif de cette partie étant avant tout de rappeler des résultats utiles pour la suite du mémoire.

\section{Introduction et postulats}

La mécanique quantique est \textbf{régie par des postulats}, l'équivalent des axiomes en mathématiques, ces postulats sont assumés pour vrai.\vspace{5pt}

On l'a vu en introduction, une particule quantique est représentée par une fonction d'onde appelée état, qu'on manipule par des opérateurs auto-adjoints qui sont appelés observables.\vspace{5pt}

On va maintenant décrire certains postulats de la mécanique quantique qui seront utiles pour expliquer et justifier l'introduction de certaines définitions ou propriétés.\vspace{4pt}


\subsection{Postulats}
La mécanique quantique affirme que les particules ont un aspect ondulatoire. Une des expériences les plus illustratrices de ce point est celle des fentes de Young. Pour les décrire, on introduit la notion de fonction d'onde. L'évolution de ces fonctions d'ondes vérifient une équation différentielle appelée équation de Schrödinger.

On identifiera par la suite les fonctions d'ondes aux états.

\begin{definition}[État pur]
    Soit $\mathcal{H}$ un espace de Hilbert, un \textbf{état pur} de $\mathcal{H}$ est une classe d"équivalence de vecteurs unitaires de $\mathcal{H}$ pour la relation d'équivalence $\psi' \simeq \psi$ si et seulement si $\psi ' = z \psi$ avec $z \in \mathbb{C}$ de module $1$.
\end{definition}

On appelle n'importe quel représentant de ces classes d'équivalence un \textbf{état}.
\begin{remark}
    On généralise en arrêtant de parler de fonctions d'ondes. Pour la suite de ce mémoire, les états seront considérés comme des éléments de $\mathcal{H} = L^2(\R^n) \bigcap \mathcal{C}^0(\R^n)$. Pour tout instant $t$, on aura un espace d'état $H_t = L^2(\R^n) \bigcap \mathcal{C}^0(\R^n)$.
\end{remark}\vspace{10pt}

Introduite par Paul Dirac en 1939, la notation bra-ket a pour but de souligner l'aspect vectoriel d'un état quantique, son nom vient de $\langle \cdot, \cdot \rangle$ prononcé bracket en anglais. Bien que cette notation ne sera pas utilisée dans la suite de ce mémoire, je trouvais intéressant de la présenter. \vspace{5pt}

\begin{definition}[Notation bra-ket]
    .
    \begin{itemize}
        \item \textbf{Les Kets} : Soit $\psi$ un vecteur de l'espace des états $\mathcal{H}$, il est souvent noté $| \psi \rangle$, ils forment un espace vectoriel complexe. On introduit maintenant leur duaux : les bras
        \item \textbf{Les bras} : L'espace duale des états est noté $\mathcal{H}^*$ et ses éléments sont appelés bras, si $v \in \mathcal{H}^*$, on l'écrit $\langle v |$
    \end{itemize}
\end{definition}

Le premier postulat est le \textbf{principe de superposition}, à chaque instant $t$, l'état d'un système est entièrement défini par un vecteur $|\psi(t)\rangle$ d'un espace de Hilbert $\mathcal{H}$. \vspace{10pt}

\begin{definition}[Observable]
    Soit $\mathcal{H}$ un Hilbert. Une \textbf{observable} est un opérateur auto-adjoint de $\mathcal{H}$.
\end{definition}
\begin{remark}
    D'après la proposition 3, son spectre est donc réel.
\end{remark}

On comprendra l'intérêt du spectre réel dans le troisième postulat. Le prochain postulat, appelé le \textbf{principe de correspondance}, dit qu'à toute propriété observable, on va associer un opérateur auto-adjoint linéaire agissant sur les vecteurs états d'un espace de Hilbert $\mathcal{H}$. Cet opérateur est nommé \textbf{observable}. On va associer à une quantité physique $A$ l'observable $\hat{A}$, par exemple à la position $x,~y,~z,~r$ on associe la multiplication par $x,~y,~z,~r$, à l'impulsion $p_x,~ p_y,~ p_z$ on associe les opérateurs de dérivations $-i \hbar \frac{\partial}{\partial x}$, $-i \hbar \frac{\partial}{\partial y}$, $-i \hbar \frac{\partial}{\partial z}$.\vspace {10pt}

\begin{remark}
    On peut définir la valeur attendue d'une observable $A$ dans l'état $\psi$ est défini par $ \langle A \rangle_{\psi} := \langle \psi, A \psi \rangle = \langle A \psi, \psi \rangle.$ Cette valeur est réelle. On peut l'interpréter comme une moyenne, une espérance. Par exemple la valeur moyenne de la position d'une particule quantique dans un état $\psi \in L^2(\R^n)$ est $\langle P \rangle _\psi$.
\end{remark}\vspace{5pt}

Le \textbf{troisième postulat} porte sur la valeur possible des observables, la mesure d'une grandeur physique représentée par l'observable $A$ ne peut fournir que l'une des valeurs propres de $A$. Cela explique que dans la définition d'une observable il y a la condition nécessaire d'être auto-adjoint.

\subsection{Interprétation probabiliste}

Les physiciens se sont demandés si l'on pouvait donner une interprétation physique à un état quantique. En 1929, Max Born a soumis l'idée que l'on pourrait interpréter le carré de l'état comme une densité de probabilité. Cette idée fait partie du courant de pensée de l'interprétation de Copenhague qui donne une interprétation cohérente de la mécanique quantique. C'est pour cela que l'on est amené à travailler avec l'espace d'états $\mathcal{H} = L^2(\R^n)$. \vspace{10pt}

\section{Transformée de Fourier}

L'objectif de cette partie est d'étudier la transformée de Fourier et de lui en donner une interprétation physique.
Comme on est amené à travaillé sur des éléments de $L^2(\R^n)$, on va devoir introduire des notations qui seront utiles afin de manipuler simplement et de façon compréhensive des points de $\R^n$.\vspace{10pt}

Les variables seront souvent notées $(\xi, x)$, ou $(p,q)$. En général dans l'espace des phase, on aura $p \in \mathbb{R}^n$ pour l'impulsion, $q \in \mathbb{R}^n$ pour la position, ces mêmes variables peuvent être notées $(\xi, x)$, dans l'espace des opérateurs pseudo-différentiels.

\subsection{Notations multi-indices}

\begin{definition}
    On introduit l'opérateur de dérivation $\partial_i = \frac{\partial}{\partial x_i}$ usuel par rapport à la i-ème variable et un opérateur nommé $D_j = -i \partial_j$. L'intérêt de ce dernier opérateur sera clarifié plus tard, il servira avant tout à simplifier les propriétés de la transformée de Fourier.
\end{definition}

On appelle multi-indice de $\mathbb{R}^n$ un $n$-uplet d'entiers $\alpha = (\alpha_1, \cdots, \alpha_n)$.\vspace{2pt}

On note alors, pour $\alpha$ un multi-indice : $\partial^\alpha = \partial_1^{\alpha_1} \cdots \partial_n^{\alpha_n}$ et $D^\alpha = D_1^{\alpha_1} \cdots D_n^{\alpha_n}$. \vspace{2pt}

On note aussi, pour $\alpha$ et  $\beta$ deux multi-indice, $\alpha! = \alpha_1! \cdots \alpha_n!$, $\binom{\beta}{\alpha} = \binom{\beta_1}{\alpha_1} \cdots \binom{\beta_n}{\alpha_n}$. Enfin on dit que $\alpha \leq \beta$ si $\alpha_j \leq \beta_j$  pour tout $j$.

\subsection{Étude de la transformée de Fourier}

Commençons par introduire des espaces utiles pour tout la suite du mémoire, puisque l'on commencera souvent par établir des résultats sur ces espaces, car ils sont plus simple à démontrer, et utiliser leur densités dans les espaces $L^p$ pour étendre les résultats aux espaces $L^p$.

\begin{definition}[Fonction à décroissance rapide (sur $\mathbb{R}^n$)]
    Une fonction $f : \mathbb{R}^n \rightarrow \mathbb{R}^n$ est dite "à décroissance rapide" si elle est indéfiniment différentiable et telle que $ \underset{x \in \mathbb{R}^n}{sup} |x^\alpha (\partial ^\beta \varphi)(x)| < \infty$ pour tout multi-indice $\alpha$ et $\beta$.
\end{definition}

\begin{definition}[Espace de Schwartz (sur $\mathbb{R}^n$)]
    L'espace de Schwartz $\mathcal{S}(\mathbb{R}^n)$ est l'ensemble des fonctions à décroissance rapide sur $\mathbb{R}^n$.
\end{definition}

\begin{definition}[Fonction $C_c^\infty(\R^n)$]
    On dit qu'une fonction de $\R^n$ appartient à $C_c^\infty(\R^n)$ si elle est lisse à support compact dans $\R^n$.
\end{definition}

\begin{theorem}
    Les espaces $C_c^\infty(\R^n)$ et $\mathcal{S}(\R^n)$ sont denses dans les espaces $L^p(\R^n)$,$ p \in [1, \infty [$.
\end{theorem}

\begin{remark}
    Ce théorème ne sera pas démontré, on admet le résultat. Il existe une démonstration dans le cours d'analyse 2.
\end{remark}

\begin{definition}[Transformée de Fourier d'une fonction $L^1(\mathbb{R}^n)$]
    On note la transformée de Fourier d'une telle fonction f par $\hat{f}$, et on a
    $$\hat{f}(\xi) := (2 \pi)^{-n/2} \int _{\mathbb{R}^n} e^{-ix \xi} f(x) dx ~~ \xi \in \mathbb{R}^n$$
\end{definition}

La transformée de Fourier fournit un lien profond entre l'étalement spatiale d'un objet, par exemple une onde, une particule, et son étalement en fréquence. Elle est ainsi très utilisée dans l'ingénierie du son, ou par exemple pour étudier la propagation de la chaleur. Ce lien est expliqué par le fait que l'on peut interpréter la transformée de Fourier en chaque point comme un produit scalaire d'une exponentielle oscillante d'une certaine fréquence et d'une fonction f, on obtient ainsi naturellement la décomposition de f en une somme infiniment continue de ses fréquences.

Pour une fonction d'onde $\psi : x \mapsto \psi(x)$ sur l'espace position on notera sa $\varphi : p \mapsto \varphi(p)$ sa transformée de Fourier sur l'espace impulsion. Quand on voit les variables $x$ et $p$ dans un couple, $x$ fera souvent référence à la position et $p$ à l'impulsion.


\begin{remark}
    $\| \psi \| = 1$ si et seulement si $\| \varphi \| =1$
\end{remark}

\begin{remark}
    On peut donner une autre écriture de la transformée de Fourier :
    $$\varphi(p) := \frac{1}{\sqrt{2 \pi \hbar}} \int_{\mathbb{R}^n} e^{-ix \frac{p}{\hbar}} \psi(x) dx$$
\end{remark}

On aura besoin du résultat suivant dans le chapitre 2, il sera admis, la démonstration utilisant des résultats d'analyse complexe.

\begin{theorem}
    Soit $a$ un réel strictement positif. Alors pour tout $\xi$ réel, la transformée de Fourier d'une gausienne $f$ définie par $f(x) = e^{-ax^2}$ est encore une gausienne définie par
    $$\hat{f}(\xi) = \sqrt{\frac{\pi}{a}}e^{\frac{- \xi ^2}{4a}}$$
\end{theorem}
\subsection{Étude de la position et de l'impulsion et principe d'incertitude de Heisenberg}

On va maintenant étudier les propriétés générales de la transformée de Fourier, il est important de bien maîtriser la transformée de Fourier et ses propriétés car elles seront souvent implicitement utilisées dans les démonstrations futures.

\begin{proposition}[Le lemme de Riemann-Lebesgue]
    Soit $f \in L^1(\R^n)$, alors $\hat{f}$ est une fonction continue sur $\mathbb{R}^n$ telle que $\underset{|\xi| \ra \infty}{lim} \hat{f}(\xi) = 0$.
\end{proposition}

\begin{proposition}
    Soit $f$ et $g$ dans $L^1(\R^n)$, alors
    $$(f*g) \hat{~} (\xi) = (2 \pi)^{n/2} \hat{f} (\xi) \hat{g} (\xi),~~ \xi \in \R^n.$$
\end{proposition}

\begin{proposition}
    Soit $\varphi \in L^1(\R^n)$. Alors pour tout multi-indices $\alpha$, \vspace{5pt}


    \begin{itemize}
        \item $(D^\alpha \varphi) \hat{~} (\xi) = \xi ^\alpha \hat{\varphi} (\xi) ,~~ \xi \in \R^n$.
        \item $(D^\alpha \hat{\varphi}) (\xi) = ((-x)^\alpha \varphi) \hat{~} (\xi),~~ x \in \R^n,~ \xi \in \R^n$.
    \end{itemize}
\end{proposition}

Pour mieux comprendre comment la transformée de Fourier fonctionne et comment la position et l'impulsion d'une particule quantique sont reliés à travers elle, on peut introduire les opérateurs suivants et regarder comment elle agit sur eux.

\begin{definition}[Opérateurs de translation, dilatation et modulation]
    On note les opérateurs suivants sur l'espace des fonctions : \vspace{5pt}
    \begin{itemize}
        \item L'opérateur de \textbf{translation} : $T_yf~(x) := f(x+y),~~~x \in \mathbb{R}^n$.
        \item L'opérateur de \textbf{modulation} : $M_yf~(x) := e^{ixy}f(x),~~~x \in \mathbb{R}^n$.
        \item L'opérateur de \textbf{dilatation} : $D_af~(x) := f(ax),~~~x \in \mathbb{R}^n$.
    \end{itemize}
    pour tout $y \in \mathbb{R}^n$ et $~a \in \mathbb{R}^*$.
\end{definition}

\begin{proposition}[Comment agit la transformée de Fourier sur les opérateurs de translation, modulation et dilatation]
    .

    \vspace{5pt}
    Soit $f \in L^1(\mathbb{R}^n)$, on a : %dans L2 ?
    \begin{itemize}
        \item $\widehat{T_yf}(\xi) = M_y\hat{f}(\xi)$.
        \item $\widehat{M_yf}(\xi) = T_{-y}\hat{f}(\xi)$.
        \item $\widehat{D_af}(\xi) = |a|^{-n}(D_{\frac{1}{a}} \hat{f})(\xi)$.
    \end{itemize}
    pour tout $\xi \in \mathbb{R}^n$.
\end{proposition}

\begin{remark}
    Le troisième point signifie que contracter une fonction dans la variable de Fourier revient à la dilater dans l'autre variable.
\end{remark}
On comprend ainsi que plus une particule quantique sera localisée dans l'espace, plus son impulsion sera très étalée, inversement si une particule quantique à une impulsion très localisée, elle sera très étalée dans l'espace des positions.


\begin{theorem}[La formule d'inversion de Fourier]
    La transformée de Fourier est une application bijective de $S(\mathbb{R}^n)$ sur $S(\mathbb{R}^n)$, de plus $\check{\hat{f}} = f, ~f \in S(\mathbb{R}^n)$ où
    $$ \check{g} = (2 \pi)^{-n/2} \int _{\mathbb{R}^n} e^{ix \xi} g( \xi) d \xi$$
\end{theorem}

Il est bon d'étendre cet opérateur aux fonctions de carrés intégrables, justifié par le postulat de la mécanique quantique de l'interprétation probabiliste des fonctions d'ondes.

\begin{theorem}[Le théorème de Plancherel]
    L'application $\mathcal{F} : S(\mathbb{R}^n) \rightarrow S(\mathbb{R}^n)$ peut être étendue de façon \textbf{unique} en un opérateur unitaire de $L^2(\mathbb{R}^n)$
\end{theorem}

\begin{proposition}[$\mathcal{F}$ est un opérateur symétrique.]
    Soit $f,~ g \in L^1(\R^n)$, alors
    $$\int_{\R^n}  \hat{f} (x) g(x) dx = \int_{\R^n} f(x) \hat{g} (x) dx.$$
\end{proposition}

On comprend ainsi comment l'espace position et impulsion sont étroitement liés. On par peut choisir d'étudier une particule dans un état quantique sur l'espace position ou l'espace impulsion, qui vont chacun décrire respectivement la position ou l'impulsion de la particule. En étudiant par exemple un électron sur l'espace position, on peut utiliser la transformée de Fourier pour étudier cet électron sur son espace impulsion.

\subsection{Introduction des opérateurs position et impulsion}

On peut maintenant introduire deux opérateurs très important en mécanique quantique, l'opérateur de position et l'opérateur d'impulsion, on aimerait représenter les coordonnées canoniques $(q,p)$ par des opérateurs auto-adjoints d'un espace de Hilbert $\mathcal{H}$, ils seront très importants pour la suite du mémoire. \vspace{10 pt}

On va les introduire par une manière naïve, celle de calculer la valeur moyenne de la position et de l'impulsion d'une particule quantique.\vspace{10 pt}

Soit une particule quantique dans l'état $\psi \in L^2(\mathbb{R}^n)$. \vspace{10 pt}

Alors \textbf{la valeur moyenne de sa position} est,
$$\langle \mathcal{Q} \rangle_\psi = \langle \psi, \mathcal{Q} \psi \rangle = \int \psi^*(x)(x\psi(x))dx$$

On définit donc l'opérateur position par
$$\mathcal{Q} : \psi(x) \mapsto x \psi(x)$$

Normalement, on s'attend à ce que l'opérateur impulsion et position soient liés par la transformé de Fourier, en effet, la transformée de Fourier relie la position et la fréquence.\vspace{5pt}

La \textbf{valeur moyenne de l'impulsion} sur l'espace impulsion est donnée par
$$ \langle \varphi, p \varphi \rangle = \int \varphi^* (p)(p \varphi(p)) dp$$

Cependant on aimerait que ces opérateurs soient définis sur le même espace position, on cherche donc à ne plus avoir une intégrale en $p$ mais une intégrale en $x$. Comme la transformée de Fourier est un opérateur symétrique, on obtient :
$$\int \psi^* (x) \frac{\hbar}{i} \frac{d \psi(x)}{dx}$$

On définit donc l'\textbf{opérateur impulsion} $\hat{p}$ par
$$\hat{P}: \psi(x) \mapsto \frac{\hbar}{i} \frac{d \psi(x)}{dx}$$

On définit finalement l'\textbf{opérateur de position} (resp. \textbf{impulsion}) par rapport à la $i$-ème(resp. $j$-ème) variable par :

$$(Q_i \psi)(x) = x_i \psi(x)$$
$$(P_j \psi)(x) =-i \hbar \frac{\partial \psi}{ \partial x_j}(x)$$

Pour tout $x \in \R^n$.
\begin{remark}
    On remarque qu'elles sont transformées de Fourier l'une de l'autre.
\end{remark} \vspace{10 pt}

Est-ce que comme en mécanique classique, $\hat{P}$ et $\hat{Q}$ commutent ? La réponse est non.

En effet, calculons le commutateur de $P_j$ et $Q_k$: \vspace{10 pt}

Soit $\psi \in L^2(\R^n)$ un état.
On a $$(Q_k \psi)(x) = x_k \psi(x) \text{ et } (P_j \psi)(x) = -i \hbar \frac{ \partial \psi}{\partial x_j}(x)$$

Donc $$Q_k(P_j \psi)(x) = -i x_k \hbar \frac{\partial \psi}{\partial x_j}(x) \text{ et } P_j(Q_k \psi)(x) = -i \hbar \frac{\partial x_k \psi}{\partial x_j}(x) = -i \hbar \left( \delta_k^j\psi(x) + x_k \frac{\partial \psi}{ \partial x_j} (x)\right)$$

Donc $$Q_k(P_j \psi)(x)- P_j(Q_k \psi)(x) = -i x_k \hbar \frac{\partial \psi}{\partial x_j}(x)-\left(-i \hbar \left( \delta_k^j \psi(x) + x_k \frac{\partial \psi}{\partial x_j}(x)\right)\right)=-i \hbar \delta_k^j \psi(x)$$

Donc $[Q_k, P_j] = \delta _k^j i \hbar Id$

\subsection{Relation de commutation canonique}
En 1925, Max Born énonce une future définition de la mécanique quantique qui affirme que deux observables qui sont conjuguées par la transformée de Fourier ( c'est à dire que l'une est transformée de Fourier de l'autre) satisfont la \textbf{relation de commutation canonique}, c'est à dire qu'elles commutent à un facteur $ i \hbar$ près. C'est bien différent de la mécanique classique où toutes les valeurs physiques commutent pour le produit de fonctions.

\begin{remark}
    On voit que le terme $i \hbar$ est vraiment très petit, ridiculement petit, et pourtant c'est ce minuscule terme qui empêche la mécanique quantique d'être égale à la mécanique classique, et qui a autant de répercussions sur cette théorie.
\end{remark}

\vspace{5pt}

\section{Opérateurs pseudo-différentiels}

On va introduire maintenant la \textbf{théorie des opérateurs pseudo-différentiels} qui sont essentiels à la compréhension de la suite du mémoire. Ils serviront essentiellement à justifier une généralisation de $(Weyl)$.

\begin{definition}[Opérateur différentiel linéaire (sur $\mathbb{R}^n$)]
    On définit un \textbf{opérateur différentiel linéaire} (sur $\mathbb{R}^n$) par $\underset{|\alpha| \leq m}{\sum} a_\alpha (x) D^\alpha, ~~ x \in \mathbb{R}^n$ et noté $P(x, D)$.
    Où $\alpha = (\alpha_1, \cdots, \alpha_n)$ est un multi-index, $m$ est un entier naturel, $a_\alpha$ est une fonction de $\R^n$, $D^\alpha = D_{\alpha_1}^{\alpha_1} \cdots D_{\alpha_n}^{\alpha_n}$ et $D_x = -i \partial _x$
\end{definition}

On va maintenant essayer d'écrire cet opérateur différemment, en utilisant son écriture de Fourier,

Pour $f \in \mathcal{S}(\R^n)$, d'après la formule d'inversion de Fourier, on peut écrire $f$ comme : $$f(x) = \check{\hat{f}}(x) =\frac{1}{(2 \pi)^{n/2}} \check{\int}_{\mathbb{R}^n} e^{-i x t} f(t) ~ dt =\frac{1}{(2 \pi)^{n/2}} \int_{\mathbb{R}^n} e^{i \xi x} \left(\int_{\mathbb{R}^n} e^{-ixt} f(t) dt\right) ~ d \xi$$

Alors $$P(x,D)(f(x)) = P(x,D)\left(\frac{1}{(2 \pi)^{n/2}}\int_{\mathbb{R}^n} e^{i \xi x} \left(\int_{\mathbb{R}^n} e^{-ixt} f(t) dt\right) ~ d \xi\right)$$
$$= \underset{|\alpha| \leq m}{\sum} a_\alpha (x) D^\alpha \left(\frac{1}{(2 \pi)^{n/2}}\int_{\mathbb{R}^n} e^{i \xi x} \left(\int_{\mathbb{R}^n} e^{-ixt} f(t) dt\right) ~ d \xi\right)$$
$$= \underset{|\alpha| \leq m}{\sum} a_\alpha (x) \left(\frac{1}{(2 \pi)^{n/2}}\int_{\mathbb{R}^n} D^\alpha (e^{i \xi x}) \left(\int_{\mathbb{R}^n} e^{-ixt} f(t) dt\right) ~ d \xi\right)$$
$$= \underset{|\alpha| \leq m}{\sum} a_\alpha (x) \left(\frac{1}{(2 \pi)^{n/2}}\int_{\mathbb{R}^n} \xi^\alpha e^{i \xi x} \left(\int_{\mathbb{R}^n} e^{-ixt} f(t) dt\right) ~ d \xi\right)$$
$$=\frac{1}{(2 \pi)^{n/2}} \int_{\mathbb{R}^n} \underset{symbole}{\underbrace{\underset{|\alpha| \leq m}{\sum} a_\alpha (x) \xi^\alpha}} e^{i \xi x} \left(\int_{\mathbb{R}^n} e^{-ixt} f(t) dt\right) ~ d\xi$$

On appelle alors $\underset{|\alpha| \leq m}{\sum} a_\alpha (x) \xi^\alpha$ le \textbf{symbole} de l'opérateur différentiel $P(x,D)$. Car formellement on remarque que $P(x,D)$ ne dépend que se ce symbole. On peut ainsi donner une définition plus large des opérateurs différentiels, en les définissant par leurs symboles et non plus par un opérateur. On les appelle opérateurs pseudo-différentiels. On sera amené à choisir des espaces de symboles plus grand dans la suite.

\begin{definition}[Symbole]
    Soit $m \in \mathbb{R}$. On définit comme $S^m$ l'ensemble des fonctions $~\mathcal{C}^\infty(\mathbb{R}^{2n})$ telles que pour tous multi-indices $\alpha$ et $\beta$, il existe une constante positive $C_{\alpha, \beta}$ dépendant uniquement de $\alpha$ et $\beta$, tel que l'on ait :
    $$|(D_x^\alpha D_\xi^\beta \sigma)(x, \xi)| \leq C_{\alpha, \beta}(1+|\xi|)^{m-|\beta|},~ x,~ \xi \in \mathbb{R}^n$$

    On appelle n'importe quel $\sigma \in \underset{m \in \mathbb{R}}{\bigcup} S^m$ \textbf{un symbole}.
\end{definition}

\begin{remark}
    On demande à contrôler la croissance par rapport à la variable $\xi$, l'intérêt est dans son utilisation future dans la transformée de Fourier afin d'assurer la convergence.
\end{remark}

On peut donc maintenant associer de façon beaucoup plus générale un symbole à un opérateur, que l'on appellera opérateur pseudo-différentiel.

\begin{definition}[Opérateur pseudo-différentiel associé à un symbole]

    On définit l'\textbf{opérateur pseudo-différentiel} $T_\sigma$ correspond au symbole $\sigma$ par
    $$(T_\sigma \varphi)(x) = (2\pi)^{-n/2} \int_{\mathbb{R}^n} e^{ix \xi} \sigma(x, \xi) \hat{\varphi} (\xi) d \xi, ~~ x \in \mathbb{R}^n$$
    pour tout $\varphi \in S(\mathbb{R}^n)$
\end{definition}

On remarque le lien avec le \textbf{problème de la quantification}, on a pu associer à une fonction particulière de $\mathbb{R}^{2n}$ un opérateur. On peut donc maintenant commencer à étudier l'opérateur proposé par Weyl $(Weyl)$ en le généralisant aux opérateurs pseudo-différentiels, formulée par Hörmander en.

\begin{proposition}
    Soit $\sigma$ un symbole. Alors l'opérateur pseudo différentiel $T_\sigma$ est une \textbf{application continue} de $\mathcal{S}(\mathbb{R}^n)$ dans $\mathcal{S}(\mathbb{R}^n)$.
\end{proposition}

\emptypage
\emptypage
\chapter{La transformée de Weyl} % au moins 10 pages

\section{La transformée de Fourier-Wigner} % 3 pages

\subsection{Théorème de Stone et groupe à 1-paramètre}

Le but de cette section est de justifier mathématiquement certains calculs proposés par la suite.

\begin{definition}[Groupe unitaire à un paramètre fortement continue]
    Une famille d'opérateurs $\{U_t\}_{t\in\R}$ sur $\mathcal{H}$ est appelée \textbf{groupe unitaire à un paramètre fortement continue} si elle satisfait les conditions suivantes :
    \begin{itemize}
        \item Pour tout $t \in \R$, $U_t : \mathcal{H} \ra \mathcal{H}$ est un opérateur unitaire.
        \item Pour tout $t, s \in \R$, $U_{t+s}=U_tU_s$. En particulier, $U_0 = id_\mathcal{H}$ et $U_{-t} = U_t^{-1} = U_t^*$
        \item Pour tout $\psi \in \mathcal{H}$, l'application $t \mapsto U_t \psi$ est continue de $\R$ dans $\mathcal{H}$.
    \end{itemize}
\end{definition}

\begin{theorem}[Théorème de Stone (forme particulière)]
    Soit $A$ un opérateur auto-adjoint sur un espace de Hilbert $\mathcal{H}$. Alors $A$ génère un groupe unitaire à un paramètre fortement continue, noté $U_t = e^{itA}$
\end{theorem}

\begin{remark}
    Ce théorème nous permet de justifier le calcul d'exponentielles d'opérateurs dans la suite du mémoire. Ce théorème ne sera pas démontré, la démonstration étant hors cadre et trop difficile.
\end{remark}

\subsection{Groupe de Heisenberg et représentation $\rho$}

Dans cette partie sera abordée brièvement la notion du \textbf{groupe de Heisenberg et de ses représentations unitaires}.

Dans toute la suite de cette partie, on identifie $(p, q) \in \R^{2n}$ à $q + ip \in \mathbb{C}^n$.

Les formes symplectique n'étant pas enseignée en M1 PMG, une définition en saura donnée.

\begin{definition}[Forme symplectique]
    Une \textbf{forme symplectique} est une application $\omega : E^2 \ra \R$ sur un espace vectoriel réel $E$, qui est bilinéaire, anti-symétrique et non-dégénérée.
\end{definition}

On définit une \textbf{forme symplectique} sur $\mathbb{C}^n$ par $[z, w] = 2 Im(z \. \bar{w}), z, w \in \mathbb{C}^n$, où $z = (z_1, \cdots, z_n)$, $w = (w_1, \cdots, w_n)$ et $z \. \bar{w} = \underset{j=1}{\overset{n}{\sum}} z_j \bar{w}_j$.

On définit une \textbf{multiplication} sur $\mathbb{C}^n \times \R$ par $(z,t) \. (w, s) = (z+w, t+s + [z,w])$ pour tout $(z,t), (w,s) \in \mathbb{C}^n \times \R$.

\begin{proposition}
    $\mathbb{C}^n \times \R$ est un \textbf{groupe} pour la multiplication définie précédemment.
\end{proposition}

Le groupe $\mathbb{C}^n \times \R$ est noté $H^n$ et est appelé \textbf{le groupe de Heisenberg}.\vspace{5pt}

\begin{remark}
    La droite réelle $\R$ est vue comme l'espace temporelle dans lequel $t \in \R$ représente un instant $t$. Et un point de $\mathbb{C}^n$ est vu comme un point dans lequel sa partie réelle élément de $\R^n$ représente par exemple sa position et sa partie imaginaire son impulsion. Il s'agit avant tout de mon interprétation personnelle.
\end{remark}

Pour tout $\lambda \in \R$, on définit $R^\lambda : H^n \ra G$, où $G$ est le groupe des opérateurs unitaires de $L^2(\R^n)$ par
$$(R^\lambda(z,t)f)(x) = e^{i \lambda (q\. x + \frac{1}{2} q \. p+ \frac{1}{4}t)}f(x+p), x \in \mathbb{R}^n$$\vspace{5pt}

Ce sont des représentations unitaires de $H^n$.\vspace{5pt}

On va s'intéresser à l'opérateur unitaire $\rho(p,q) := R^1(z,0)$ où $z=q+ip$. \vspace{10pt}

On va maintenant motiver les objets précédents et justifier les futures calculs de la transformée de Weyl avec les résultats suivant. A partir de maintenant $\hat{X}$, $\hat{Y}$ et $\hat{Z}$ désignent des opérateurs. On utilise le commutateur usuel de deux opérateurs $\hat{X}$ et $\hat{Y}$ définis par $[\hat{X},\hat{Y}] = \hat{X}\hat{Y}-\hat{Y}\hat{X}$. On considère que c'est un crochet de Lie avec les propriétés qui en découlent.\vspace{5pt}

On introduit la formule de \textbf{Baker-Campbell-Hausdorff} qui donne la solution $\hat{Z}$ de l'équation $e^{\hat{Z}} = e^{\hat{X}} e^{\hat{Y}}$. On a alors l'élégante formule :
$$e^{\hat{A}} e^{\hat{B}} = e^{\hat{A}+\hat{B}+\frac{1}{2} [\hat{A},\hat{B}] + \frac{1}{12} ([\hat{A}, [\hat{A},\hat{B}]] + [\hat{B},[\hat{B},\hat{A}]])+\cdots}$$

Si $\hat{A}$ et $\hat{B}$ commutent avec leur commutateurs c'est-à-dire $[\hat{A},[\hat{A}, \hat{B}]] = 0$ et $[\hat{B}, [\hat{B}, \hat{A}]] = 0$.

Cette formule devient la \textbf{formule dite de Glauber} :
$$e^{\hat{X}+\hat{Y}} = e^{\hat{X}} e^{\hat{Y}}e^{-\frac{1}{2}[\hat{X}, \hat{Y}]}$$\vspace{5pt}

La formule de Baker-Campbell-Hausdorff et a été introduite dans ce mémoire pour avoir une compréhension plus fine sur les exponentielles d'opérateurs, c'est la formule de Glauber qui sera utilisée dans la suite.\vspace{5pt}

On va calculer des résultats annexes utilisés dans le début de la partie sur la transformation de Weyl. Ils ne seront pas justifiés mais pourront fournir une motivation suffisante à la généralisation de la formule $(Weyl)$ en un opérateur pseudo-différentiel.\vspace{5pt}


On va calculer l'exponentielle de l'opérateur de position et d'impulsion respectivement, on peut en donner une brève justification avec le théorème de Stone et les groupes à 1-paramètres fortement continus.

\textbf{Calcul de $e^{\xi \mathcal{\hat{Q}}}$}

Soit $f \in \mathcal{S}(\R^n),~ \xi \in \R^n$. Alors,

$$(e^{\xi \mathcal{\hat{Q}}}f)(x) = \left(\underset{k=0}{\overset{+ \infty}{\sum}} \frac{(\xi \mathcal{\hat{Q}})^k}{k!}f\right)(x) = \underset{k=0}{\overset{+ \infty}{\sum}} \frac{(\xi x)^k f(x)}{k!} = f(x) \underset{k=0}{\overset{+ \infty}{\sum}} \frac{(\xi x)^k}{k!} = f(x) e^{\xi x} = Id ~e^{\xi x}$$

Avec $Id$ qui représente l'opérateur identité et $x$ l'opérateur multiplication par $x$.\vspace{10pt}

\textbf{Calcul de $e^{p \mathcal{\hat{P}}}$}
Soit $f \in \mathcal{S}(\R^n),~ p \in \R^n$. Alors,

$$(e^{p \mathcal{\hat{P}}}f)(x) = \left(\underset{k=0}{\overset{+ \infty}{\sum}} \frac{(p \mathcal{\hat{P}})^k}{k!}f\right)(x) = \underset{k=0}{\overset{+ \infty}{\sum}} \frac{p^k f^{(k)}(x)}{k!} = f(p+x) = T_p(x)$$

On reconnaît la série entière de $f$,

$$= \underset{k=0}{\overset{+ \infty}{\sum}} \frac{f^{(k)}(x)}{k!}((p+x)-x)^k = (T_pf)(x)$$

où $T_p$ est l'opérateur de translation par $p$ qui translate $f$ par $p$.

\begin{remark}
    On peut voir l'ensemble des opérateurs de translation comme un groupe à 1-paramètre fortement continue.
\end{remark}

\subsection{La transformée de Fourier-Wigner}

Cette partie aura pour but de présenter la transformée de Fourier-Wigner.

Soit $q$ et $p$ dans $\R^n$, et f une fonction mesurable sur $\R^n$. On définit la fonction $\rho (q,p) f$ sur $\R^n$ par
$$(\rho(q,p)f)(x) = e^{iq\.x + \frac{1}{2} iq \.p}f(x+p), ~~x\in\R^n.$$

\begin{proposition}
    $\rho(q,p)f : L^2(\R^n) \ra L^2(\R^n)$ est un \textbf{opérateur unitaire} pour tout $q$ et $p$ dans $\R^n$.
\end{proposition}


\begin{definition}
    Soit $f$ et $g$ dans $\mathcal{S}(\R^n)$, on définit alors la fonction $V(f,g)$ sur $\R^{2n}$ par
    $$V(f,g)(q,p) = (2 \pi)^{-n/2} \langle \rho (q,p)f,g \rangle, ~~q,p \in \R^n$$
    où $\langle \., \. \rangle$ est le produit scalaire de $L^2(\R^n)$.
    On appelle alors $V(f,g)$ la transformée de Fourier-Wigner de $f$ et $g$.
\end{definition}

\begin{remark}
    Pour $f, ~ g \in L^2(\R^n)$, $V(f,g)$ correspond au coefficient matriciel en $(f,g)$ de la représentation $\rho$. On obtient ainsi une meilleure intuition de ce qu'elle signifie.
\end{remark}
\begin{remark}
    On peut mieux comprendre cette écriture en essayant d'écrire la transformée de Fourier sous la forme d'un produit scalaire !
    On a alors $\hat{f} (\xi) = (2 \pi)^{-n/2} \langle \rho(-\xi,0)f, Id \rangle$
\end{remark}

\begin{remark}
    C'est une transformée de Fourier tordue.
\end{remark}\vspace{5pt}

On peut alors donner une formule plus explicite de $V(f,g)$.

\begin{proposition}
    Soit $f$ et $g$ dans $\mathcal{S}(\R^n)$, alors
    $$V(f,g)(q,p) = (2 \pi)^{-n/2} \int_{R^n} e^{iq \. y} f\left(y + \frac{p}{2}\right) \overline{g\left(y-\frac{p}{2}\right)}dy$$
    pour tout q et p dans $\R^n$.
\end{proposition}

\begin{preuve}
    Soit $f$ et $g$ dans $\mathcal{S}(\R^n)$, alors d'après la définition, on a
    $$V(f,g)(q,p) = (2 \pi)^{-n/2} \langle \rho(q,p)f,g \rangle$$
    Comme $\rho(p,q)$ est un \textbf{opérateur unitaire}, on peut l'écrire sous la forme d'une intégrale, la convergence de l'intégrale étant justifiée par l'\textbf{inégalité de Cauchy-Schwartz}.

    $$(2 \pi)^{-n/2} \int_{\R^n}e^{iqx+\frac{1}{2}iqp}f(x+p)\overline{g(x)}dx$$
    On fait alors le changement de variable $x = y-\frac{p}{2}$ et on obtient

    $$V(f,g)(q,p) = (2 \pi)^{-n/2} \int_{R^n} e^{iq \. y} f\left(y + \frac{p}{2}\right) \overline{g\left(y-\frac{p}{2}\right)}dy$$

\end{preuve}

Une propriété intéressante est la bilinéarité de la transformée de Fourier-Wigner.
\begin{proposition}
    $V : \mathcal{S} (\R^n) \times \mathcal{S} (\R^n) \ra \mathcal{S} (\R^{2n})$ est une \textbf{application bilinéaire}.
\end{proposition}

\subsubsection{La théorie des radars}
Une application intéressante et concrète de la transformée de Fourier-Wigner est dans la théorie des radars, un bref aperçu sera donné, sans rentrer dans trop de détails techniques.
Plaçons nous dans le cas d'un appareil radar qui transmet un signal électromagnétique qui rebondi sur la cible et revient à l'appareil. On peut représenter le signal par une fonction à valeurs complexes $f(t) = f_0(t) e^{2 \pi i w t}$, où $w$ représente la fréquence et $t$ le temps. Le signal revenu n'est pas le même signal que celui émis, disons que, après quelques approximations, le signal reçu est $f_{\tau \phi}(t) = f(t - \tau) e^{2 \pi i \phi t}$, où $\tau$ représente le délai de retour du signal et $\phi$ la nouvelle fréquence.
On se place maintenant dans la situation de deux cibles qui produisent des signaux $f_{\tau_1 \phi_1}$ et $f_{\tau_2 \phi_2}$. Le problème est que si ces signaux se ressemblent, il sera difficile de distinguer les 2 cibles.
On peut alors s'intéresser à la moyenne quadratique de la différence.
$$\int |f_{\tau_1 \phi_1} - f_{\tau _2 \phi_2}|^2(t) ~dt$$
En développant la valeur absolue au carré, on obtient
$$2 \int |f|^2 dt -2 Re \langle f_{\tau_1 \phi_1} , f_{\tau _2 \phi_2} \rangle$$
On va alors étudier le second terme qui ne dépend que des variables. \vspace{3pt}

On a $$\langle f_{\tau_1 \phi_1} , f_{\tau _2 \phi_2} \rangle = e^{2 \pi i w (\tau _1 - \tau _1)} \int f_0(t- \tau_1) \overline{f_0(t-\tau_2)}e^{2 \pi i(\phi_1 - \phi_2)t}dt$$

On voit que ce terme ressemble à transformée de Fourier-Wigner, en faisant un changement de variable adapté on trouve que ce terme peut s'exprimer comme la transformée de Fourier-Wigner d'une certaine fonction appelée fonction d'ambiguïté qui est égale à $V(f,f)$.


\section{La transformée de Wigner} % 3 pages

On introduit maintenant la transformée de Wigner et certaines de ses propriétés. On peut considérer la transformée de Wigner $W(f)$ d'une fonction $f \in L^2(\R^n)$.% pour étudier "the non existing joint probability distribution of position and momentum in the state f".

On commence par calculer la transformée de Fourier de la transformée de Fourier-Wigner.

\begin{theorem}
    Soit $f$ et $g$ dans $\mathcal{S}(\R^n)$. Alors
    $$V(f,g)\hat{~}(x, \xi) = (2 \pi) ^{-n/2} \int_{\R^n} e^{-i \xi p} f\left(x+ \frac{p}{2}\right) \overline{g\left(x-\frac{p}{2}\right)} dp, x, \xi \in \R^n$$
\end{theorem}

\begin{remark}
    On note que cette expression ne diffère de la transformée de Fourier-Wigner que d'un signe $-$ dans le coefficient de l'exponentielle.
\end{remark}

\begin{preuve}
    Soit $\epsilon >0$, on définit la fonction $I_\epsilon$ sur $\R^{2n}$ par
    $$I_\epsilon(x, \xi) = \int_{\R^n} \int_{\R^n} e^{-\frac{\epsilon ^2 |q|^2}{2}} e^{-ix \. q - i \xi \. p}V(f,g)(q,p)~dqdp,~ x,~ \xi \in \R^n$$

    On a alors, en décomposant $V(f,g)$,
    $$I_\epsilon(x, \xi) = (2 \pi)^{-n/2}\int_{\R^n} \int_{\R^n} e^{-\frac{\epsilon ^2 |q|^2}{2}} e^{-ix \. q - i \xi \. p} \times \left(\int_{\R^n} e^{iq\.y} f\left(y+\frac{p}{2}\right) \overline{g\left(y-\frac{p}{2}\right)}dy\right) dp dq$$

    Par Fubini, on a :
    $$= (2 \pi)^{-n/2}\int_{\R^n} e^{ - i \xi \. p} \left(\int_{\R^n}\left(\int_{\R^n}  e^{-i(x-y)\.q}e^{-\frac{\epsilon ^2 |q|^2}{2}} dq\right) f\left(y+\frac{p}{2}\right) \overline{g\left(y-\frac{p}{2}\right)} \right) dy dp$$

    On essaye de faire apparaître autant que possible $V(f,g)$,
    $$= (2 \pi)^{-n/2}\int_{\R^n} e^{ - i \xi \. p} \left(\int_{\R^n} \epsilon^{-n} e^{-\frac{|x-y|^2}{2 \epsilon^2}} f\left(y+\frac{p}{2}\right) \overline{g\left(y-\frac{p}{2}\right)} dy \right)dp$$

    Maintenant, pour tout $p \in \R^n$, on définit la fonction $F_p$ sur $\R^n$ par
    $$F_p(y) = f\left(y+\frac{p}{2}\right)\overline{g\left(y-\frac{p}{2}\right)},~ y \in \R^n$$

    Et on définit la fonction $\varphi_\epsilon$ par $\varphi_\epsilon(y) = \epsilon^{-n} e^{-\frac{|x-y|^2}{2 \epsilon ^2}}$\vspace{5pt}

    Alors on a,
    $$I_\epsilon (x, \xi) = \int_{\R^n} e^{-i \xi \. p}(F_p * \varphi_\epsilon)(x)dp$$

    Pour $p$ fixé on a,
    $$F_p * \varphi_\epsilon \ra (2 \pi) ^{n/2} F_p \text{ uniformément sur chaque compact de } \R^n \text{ quand } \epsilon \ra 0$$

    On va essayer de majorer la valeur absolue de ce qu'on intègre pour pouvoir applique le théorème de convergence dominé de Lebesgue,\vspace{5pt}

    On peut appliquer l'inégalité de Young pour la convolution avec $p =1, q = 0$ dans ce cas précis.

    $$|(F_p * \varphi _\epsilon)(x)| \leq \|F_p\|_{L^\infty(\R^n)} \|\varphi _\epsilon\|_{L^1(\R^n)}$$
    $$= \|F_p\|_{L^\infty(\R^n)} \|\varphi\|_{L^1(\R^n)}$$
    $$\leq (2 \pi) ^{n/2} \underset{y \in \R^n}{sup}\left|f(y+\frac{p}{2}) f(y- \frac{p}{2})\right|$$
    Comme ce sont des fonctions Schwarz, il existe un entier positif $N$ et une constante $C_N$ tel que
    $$\leq C_N (1+|p|^2)^{-N}, x, p \in \R^n$$
    pour tout $\epsilon \geq 0$.

    Donc on peut utiliser le théorème de convergence dominé de Lebesgue,
    $$\underset{\epsilon \ra 0}{lim} ~I_\epsilon (x \xi) = (2 \pi)^{n/2} \int_{\R^n} e^{-i \xi \.p}f(x+ \frac{p}{2}) \overline{g(x-\frac{p}{2})} dp, x, \xi \in \R^n$$

    Ensuite en réutilisant le théorème de convergence dominé de Lebesgue on montre l'autre limite.
\end{preuve}

\begin{remark}
    Cette démonstration est intéressante car elle montre une technique usuelle de démonstration, on pose un $\epsilon$ positif et on étudie un objet très proche, pour lequel les convergences et les calculs se dérouleront plus facilement, puis une fois le calcul fini on fait tendre $\epsilon$ vers $0$.

    On a introduit le terme gaussien car sa transformée de Fourier sera gausienne et donc le calcul est simplifié.
\end{remark}



On peut maintenant définir la transformée de Wigner de deux fonctions de $\mathcal{S}(\R^n)$.

\begin{definition}[Transformée de Wigner]
    Soit $f$ et $g$ dans $\mathcal{S}(\R^n)$, alors la fonction $W(f,g)$ sur $\R^{2n}$ définie par,
    $$W(f,g)(x, \xi) = (2 \pi)^{-n/2} \int_{\R^n} e^{-i \xi \; p} f\left(x+\frac{p}{2}\right) \overline{g\left(x-\frac{p}{2}\right)} dp, ~x, ~\xi \in \R^n$$,
    est appelée la \textbf{transformée de Wigner} de f et g.
\end{definition}

\begin{theorem}[Identité de Moyal]
    Pour tout $f_1, g_1, f_2, g_2 \in \mathcal{S}(\R^n)$, on a,
    $$\langle W(f_1, g_1), W(f_2, g_2) \rangle = \langle f_1, f_1 \rangle \overline{\langle g_1,g_2 \rangle}.$$
\end{theorem}

\begin{remark}
    Elle est aussi vraie pour la transformée de Fourier-Wigner. Elle ne sera pas démontrée, n'étant pas au centre de ce mémoire.
\end{remark}

Encore une fois, il est bon de chercher à étendre l'opérateur $W$ sur l'espace $L^2(\R^n)$. Justifié par l'interprétation probabiliste des états.

\begin{proposition}%c'est un corollaire
    $W : \mathcal{S}(\R^n) \times \mathcal{S}(\R^n) \ra \mathcal{S}(\R^{2n})$ peut être étendu de façon \textbf{unique} en un opérateur bilinéaire
    $$W : L^2(\R^n) \times L^2(\R^n) \ra L^2(\R^{2n})$$

    tel que
    $$\|W(f,g)\|_{L^2(\R^{2n})}=\|f\|_{L^2(\R^n)}\|g\|_{L^2(\R^n)}.$$
\end{proposition}

\begin{remark}
    Il est est de même de la transformée de Fourier-Wigner V.
\end{remark}

\begin{proposition}
    Pour tout $f, g \in \mathcal{S}(\R^n)$, on a $W(g,f) = \overline{W(f,g)}$.
\end{proposition}

\begin{proposition}%c'est un corollaire
    On pose $W(f) = W(f,f), f \in L^2(\R^n)$, alors, \vspace{3pt}
    \begin{itemize}
        \item $W(\rho(a,b)f)(x, \xi) = W(f)(x+b, \xi-a)$, $a$, $b$, $x$, $\xi \in \R^n$
        \item $W(f)$ est à valeur réelle.
    \end{itemize}
\end{proposition}

%lien avec des probas ??


\section{La transformée de Weyl} % 4 pages

On va donc maintenant introduire la transformée de Weyl et son lien avec la transformée de Wigner.

On arrive maintenant au c\oe ur de ce mémoire, mais avant d'introduire la transformée de Weyl, on va introduire des outils nécessaires afin de comprendre et motiver son introduction.\vspace{5pt}

Soit $\sigma \in S^m$, $m \in \R$, alors pour toute fonction $\varphi \in \mathcal{S} (\R^n)$, on se souvient à la fin de la première partie que
$$(T_\sigma \varphi)(x) = (2 \pi) ^{-n} \int_{\R^n} \int_{\R^n} e^{i(x-y) \. \xi} \sigma (x, \xi) \varphi(y) dy d\xi, x \in \R ^n$$\vspace{5pt}

On a une autre façon d'associer $\sigma$ à un opérateur, la formule de Hörmander,

$$(W_\sigma \varphi)(x) = (2 \pi) ^{-n} \int_{\R^n} \int_{\R^n} e^{i(x-y) \. \xi} \sigma \left(\frac{x+y}{2}, \xi\right) \varphi(y) dy d\xi,~ x \in \R ^n$$

On remarque qu'on a juste symétrisé la première variable de la fonction $\sigma$, mais ce petit changement aura des grandes conséquences. C'est un opérateur pseudo-différentiel, mais comment en est on arrivé à cette formule ? Essayons de transformer un peu la formule $(Weyl)$,

$$\sigma(\hat{P},\hat{Q}) = \frac{1}{2 \pi \hbar} \iint \hat{\sigma} (p,q) e^{ \frac{i}{\hbar}(p\hat{P}+q\hat{Q})} dpdq$$

En utilisant la formule de Glauber car $\hat{P}$ et $\hat{Q}$ commutent avec leur commutateur, et d'après les calculs de l'exponentielle des opérateurs positions et impulsion, à la fin de la partie sur les représentations du groupe de Heisenberg, on a, en développant $\hat{\sigma}$ en fonction de $\sigma$,

$$= \frac{1}{2 \pi \hbar} \iiiint \sigma (\xi,\omega) e^{-2 \pi i (p \xi + q \omega)} e^{2 \pi i q x + \pi i p q} f(x+p) d \xi d \omega dpdq$$

En intégrant par rapport à $q$, on a
$$= \frac{1}{2 \pi \hbar} \iiint \sigma (\xi,\omega) e^{-2 \pi i p \xi} f(x+p)~ \delta\left(x-\omega + \frac{p}{2}\right) d \xi d \omega dp$$
où $\delta$ est le Dirac en ce point.
$$= \frac{1}{2 \pi \hbar} \iint \sigma \left(\xi,x+\frac{p}{2}\right) e^{-2 \pi i p \xi} f(x+p)  dp d \xi $$

On pose $y=x+p$, on obtient alors,
$$= \frac{1}{2 \pi \hbar} \iint \sigma \left(\xi,\frac{x+y}{2}\right) e^{-2 \pi i (x-y) \xi} f(y) dp d \xi $$
On reconnait alors la formule plus haut à un facteur près.

\begin{remark}
    Cette définition de $\sigma(\hat{P}, \hat{Q})$ ne dépend pas de la paramétrisation de la représentation $\rho$.
\end{remark}

\vspace{10pt}


La proposition suivante est un outil pour les démonstrations futures, il ne sera pas démontré la démonstration étant trop difficile pour moi.

\begin{proposition}%lemme
    La limite
    $$\underset{\epsilon \ra 0_+}{lim} (2 \pi)^{-n} \int_{\R^n} \int_{\R^n} \theta( \epsilon \xi) e^{i(x-y) \. \xi} \sigma(\frac{x+y}{2}, \xi) \varphi(y) dy d\xi$$

    existe et est indépendante du choix de la fonction $\theta$, de plus la convergence est unifrme pour $x \in \R$
\end{proposition}

Aucune démonstration ne sera donné, on l'assumera pour vrai.


\begin{proposition}
    Soit $\sigma \in S^m, m \in \R$, alors $W_\sigma : S(\R^n) \ra S(\R^n)$ est continue.
\end{proposition}

On appelle $W_\sigma$ la transformée de Weyl associée au symbole $\sigma$. Le théorème suivant \textbf{illustre le lien} de la transformée de Weyl avec la transformée de Wigner.\vspace{5pt}

\begin{theorem}
    Soit $\sigma \in S^m, m \in \R$, alors
    $$\langle W_\sigma f, g \rangle = (2 \pi)^{-n/2} \int_{\R^n} \int_{\R^n} \sigma(x, \xi) W(f,g) (x, \xi) dx d\xi,~ f,~ g \in \mathcal{S} (\R^n)$$
\end{theorem}
\begin{preuve}
    Soit $\theta \in C_0^\infty(\R^n)$ tel que $\theta(0)=1$, alors,
    $$\int_{\R^n} \int_{\R^n} \sigma(x, \xi) W(f,g)(x,\xi) dx d\xi$$

    On introduit $\theta$,
    $$\underset{\epsilon \ra 0_+}{lim} \int_{\R^n}  \int_{\R^n} \theta( \epsilon \xi) \sigma (x, \xi) W(f,g)(x, \xi) dx d\xi$$
    $$=\underset{\epsilon \ra 0_+}{lim} (2 \pi)^{-n/2} \int_{\R^n}  \int_{\R^n}\theta( \epsilon \xi) \sigma (x, \xi) \times \left( \int_{\R^n} e^{-i \xi \. p} f\left(x+\frac{p}{2}\right) \overline{g\left(x-\frac{p}{2}\right)} dp\right) dx d\xi $$
    $$ = \underset{\epsilon \ra 0_+}{lim} (2 \pi)^{-n/2}  \theta(\epsilon \xi) \times \left( \int_{\R^n} \int_{\R^n} \sigma(x, \xi) e^{-i \xi \. p} f\left(x+\frac{p}{2}\right)\overline{g\left(x+\frac{p}{2}\right)} dp dx\right) d\xi$$

    On pose $u = x + \frac{p}{2}, v = x-\frac{p}{2}$ alors
    $$\int_{\R^n} \int_{\R^n} \sigma (x, \xi) W(f,g)(x, \xi) dx d\xi$$
    $$=\underset{\epsilon \ra 0_+}{lim} \int_{\R^n} \theta(x \xi) \left( \int_{\R^n}\int_{\R^n} \sigma\left(\frac{u+v}{2}, \xi\right) e^{i(v-u) \. \xi} f(u) \overline{g(v)} du dv \right) d \xi$$
    On échange les intégrales par le théorème de Fubini,
    $$=\underset{\epsilon \ra 0_+}{lim} \int_{\R^n} \overline{g(v)} \left( \theta(x \xi) \sigma\left(\frac{u+v}{2}, \xi\right) e^{i(v-u) \. \xi} f(u) du d\xi\right) dv$$
    $$= (2 \pi)^{-n/2} \int_{\R^n} \overline{g(v)} (W_\sigma f)(v) dv = (2 \pi)^{n/2} \langle W_\sigma f, g\rangle$$
\end{preuve}

Maintenant que l'on a vu et démontré le lien de la transformée de Weyl avec la transformée de Wigner, on va chercher à l'étendre sur l'espace $L^2(\R^n)$, et on va chercher à étendre son symbole à une classe plus large sur $L^2(\R^n)$.\vspace{7pt}

On introduit la définition d'algèbre.

\begin{definition}
    Une \textbf{algèbre} sur un corps commutatif $\mathbb{K}$, ou simplement $\mathbb{K}-$algèbre, est une structure algèbrique $(A,+, \., \times)$ telle que:
    \begin{itemize}
        \item $(A,+, \.)$ est un espace vectoriel sur $\mathbb{K}$.
        \item la loi $\times$ est définie de $A \times A$ dans A et est bilinéaire.
    \end{itemize}
\end{definition}

Une $\mathbb{C}^*$-algèbre est une algèbre de Banach involutive, c'est à dire que l'espace vectoriel sur le corps des complexes est normé et qu'il est muni d'une involution. Une involution est par exemple l'opération de prendre l'opérateur adjoint. Cet espace est donc stable par cette opération.

On appelle la $\mathbb{C}^*$-algèbre de tous les opérateurs linéaires bornés de $L^2(\R^n)$ dans $L^2(\R^n)$ par $B(L^2(\R^n))$. Les opérateurs bornés sont préservés par l'involution.

Il est bon de chercher à étendre cet opérateur.
\begin{theorem}
    Il existe un unique opérateur linéaire borné $Q : L^2(\R^{2n}) \ra B(L^2(\R^n))$ tel que
    $$\langle ( Q_\sigma)f,g\rangle = (2 \pi)^{-n/2} \int_{\R^n} \int_{\R^n} \sigma(x, \xi) W(f,g)(x, \xi) dx d\xi$$
    et $$\|Q \sigma\|_* \leq (2 \pi)^{-n/2}\|\sigma\|_{L^2(\R^{2n})}$$
\end{theorem}

\begin{preuve}
    Soit $\sigma \in \mathcal{S}(\R^{2n})$. Alors pour tout $f \in \mathcal{S}(\R^n)$, on définit $(\mathcal{Q}_\sigma) f$ par
    $$(Q_\sigma) f = W_\sigma f$$

    On cherche maintenant à \textbf{étendre} l'opérateur $\mathcal{Q_\sigma}$\vspace{5pt}

    Alors on a $$\langle(Q_\sigma)f, g \rangle = \langle W_\sigma f, g \rangle = (2 \pi)^{-n/2} \int_{\R^n} \int_{\R^n} \sigma (x, \xi) W(f,g)(x, \xi) dx d\xi$$ pour tout $f, g \in \mathcal{S}(\R^n)$.



    $$|\langle (Q_\sigma)f, g \rangle| \leq (2 \pi)^{-n/2} \|\sigma\|_{L^2(\R^n)} \|W(f,g)\|_{L^2(\R^n)} = (2 \pi)^{-n/2} \|\sigma\|_{L^2(\R^n)} \|f\|_{L^2(\R^n)} \|g\|_{L^2(\R^n)}$$ pour tout $f, g \in \mathcal{S}(\R^n)$.\vspace{5pt}

    Donc
    $$\|Q_\sigma f\|_* \leq (2 \pi)^{-n/2} \|\sigma\|_{L^2(\R^n)} \|f\|_{L^2(\R^n)}$$ pour tout $f \in S(\R^n)$, donc
    $$\|Q_\sigma\|_* \leq (2 \pi)^{-n/2} \|\sigma\|_{L^2(\R^n)}. $$\vspace{10pt}



    Maintenant soit $\sigma \in L^2(\R^n).$ Soit $\{\sigma_k\}_{k=1}^\infty$ une suite de fonctions de $\mathcal{S}(\R^{2n})$ telle que $\sigma_k \ra \sigma$ dans $L^2(\R^n)$ quand $k \ra \infty$, alors
    $$\|Q_{\sigma_k}-Q_{\sigma_l}\|_* \leq (2 \pi)^{-n/2} \|\sigma_k - \sigma _l\|_{L^2(\R^{2n})} \ra 0$$ quand $k, l \ra \infty$.\vspace{4pt}

    Alors $\{Q_{\sigma_k}\}_{k=1}^\infty$ est une suite de Cauchy dans $B(L^2(\R^n))$. Elle converge donc, comme c'est un espace de Hilbert. On définit $Q_\sigma$ sa limite dans $B(L^2(\R^n))$. \vspace{4pt}

    Cette définition est indépendante du choix de la suite, mais soit $\{\tau_k\}$ une autre suite de fonction dans $\mathcal{S}(\R^{2n})$ tel que $\tau_k \ra \sigma$ dans $L^2(\R^{2n})$ quand $k \ra \infty$. Alors,
    $$\|Q_{\sigma_k}- Q_{\tau_k}\|_* \leq (2 \pi)^{-n/2} \|\sigma _k - \tau _k\|_{L^2(\R^{2n})} \ra 0$$ quand $k \ra \infty$. En conséquence, les limites dans $B(L^2(\R^n))$ de $\{\sigma _k\}_{k=1}^\infty$ et $\{Q_{\tau_ k}\}_{k=1}^\infty$ sont égales.  \vspace{5pt}

    Maintenant, soit $\sigma \in L^2(\R^{2n})$ et soit $\sigma _k$ une suite de fonctions dans $\mathcal{S}(\R^{2n})$ telle que $\sigma _k \ra \sigma$ dans $L^2(\R^{2n})$ quand $k \ra \infty$.\vspace{5pt}

    Alors,
    $$\|Q_\sigma\|_* = \underset{k \ra \infty}{lim} \|Q_{\sigma_k}\|_* \leq  (2 \pi)^{-n/2}  \underset{k \ra \infty}{lim} \| \sigma_k \|_{L^2(\R^n)} = (2 \pi)^{-n/2} \|\sigma\|_{L^2(\R^{2n}}$$ et l'inégalité du théorème est prouvée.

    On a donc étendu le symbole, cherchons maintenant à \textbf{étendre les fonctions}.

    Maintenant, si $f$ et $g$ sont dans $\mathcal{S}(\R^n)$, alors,

    $$\langle ( Q_ \sigma)f, g \rangle = \underset{k \ra \infty}{lim} \langle  Q_{\sigma_k}f, g \rangle = \underset{k \ra \infty}{lim} (2 \pi)^{-n/2} \int_{\R^n} \int_{\R^n} \sigma_k (x, \xi) W(f,g) (x, \xi) dx d\xi$$
    $$=(2 \pi)^{-n/2} \int_{\R^n} \int_{\R^n} \sigma (x, \xi) W(f,g) (x, \xi) dx d\xi.$$

    On étend maintenant aux fonction $L^2(\R^n)$, soit $f$ et $g$ dans $L^2(\R^n)$, alors on prend les suites $\{f_k\}_{k=1}^\infty$ et $\{g_k\}_{k=1}^\infty$ dans $\mathcal{S}(\R^n)$ telle que $f_k \ra f$ dans $L^2(\R^n)$ et $g_k \ra g$ dans $L^2(\R^n)$ quand $\ra \infty$.

    On a par densité des fonctions Schwarz et parc continuité du produit scalaire,
    $$\langle (Q_\sigma)f,g \rangle = \underset{k \ra \infty}{lim} \langle (Q _\sigma)f_k, g_k \rangle = \underset{k \ra \infty}{lim} (2 \pi)^{-n/2} \int_{\R^n} \int_{\R^n} \sigma(x, \xi) W(f_k, g_k)(x, \xi) dx d\xi.$$

    On peut passer à la limite, par continuité de $W$.
    $$=(2 \pi)^{-n/2} \int_{\R^n} \int_{\R^n} \sigma(x, \xi) W(f, g)(x, \xi) dx d\xi.$$

    Il y a bien unicité de $\mathcal{Q}$.
\end{preuve}

Pour éviter les confusions, on nommera l'opérateur $\mathcal{Q}_\sigma$ par $W_\sigma$.

\subsection{La transformée de Weyl et quelques exemples de correspondance}

On va maintenant illustrer par quelques exemples la transformée de Weyl de certains symboles, et vérifier qu'elle transforme $p$ en $\hat{P}$ et $q$ en $\hat{Q}$

\subsubsection{Exemples}
Soit $x,~y \in \R^n$. Soit $\varphi \in \mathcal{S}(\R^n)$.

\begin{exemple}[Avec $\sigma (x, \xi) := x$.]
    On a
    $$(W_\sigma \varphi)(x)=(2 \pi)^{-n} \int_{\R^n} \int_{\R^n} e^{i(x-y) \. \xi}~ \sigma\left(\frac{x+y}{2}, ~\xi\right) ~ \varphi(y) ~dy d\xi$$
    $$= (2 \pi)^{-n} \int_{\R^n} \int_{\R^n} e^{i(x-y) \. \xi} ~ \frac{x+y}{2} ~ \varphi(y) ~dy d\xi$$
    $$= (2 \pi)^{-n} \int_{\R^n} \int_{\R^n} e^{i(x-y) \. \xi} ~ \frac{x}{2} ~ \varphi(y) ~dy d\xi + (2 \pi)^{-n} \int_{\R^n} \int_{\R^n} e^{i(x-y) \. \xi} ~ \frac{y}{2} ~ \varphi(y) ~dy d\xi$$
    $$= (2 \pi)^{-n} ~ \frac{x}{2} ~ \int_{\R^n} e^{ix \. \xi} \left( \int_{\R^n} e^{-iy \. \xi} \. \varphi(y) ~dy\right) d\xi + (2 \pi)^{-n} ~ \frac{1}{2} ~\int_{\R^n} e^{ix \. \xi} \left(\int_{\R^n} e^{-iy \. \xi} \. y \. \varphi(y) ~dy \right) d\xi$$

    $$= (2 \pi)^{-n} ~ \frac{x}{2} ~ \int_{\R^n} e^{ix \. \xi} ~ \hat{\varphi}(\xi) ~ d\xi - (2 \pi)^{-n} ~ \frac{1}{2} ~\int_{\R^n} e^{ix \. \xi} \left(\int_{\R^n} D_\xi( e^{-iy \. \xi}) \. \varphi(y) ~dy \right) d\xi$$
    $$= \frac{x}{2} \varphi(x) - (2 \pi)^{-n} ~ \frac{1}{2} ~\int_{\R^n} e^{ix \. \xi}~ (D_\xi\hat{\varphi})(\xi) ~ d\xi$$
    $$= \frac{x}{2} \varphi(x) - (2 \pi)^{-n} ~ \frac{1}{2} ~\int_{\R^n} e^{ix \. \xi}~ (-x \varphi)\hat{~}(\xi) ~ d\xi$$
    $$= \frac{x}{2} \varphi(x) + \frac{x}{2}\varphi(x) = x \varphi(x) = (\mathcal{Q} \varphi)(x)$$
\end{exemple}

\begin{exemple}[Avec $\sigma (x, \xi) := \xi$.]
    On a
    $$(W_\sigma \varphi)(x)=(2 \pi)^{-n} \int_{\R^n} \int_{\R^n} e^{i(x-y) \. \xi}~ \sigma(\frac{x+y}{2}, ~\xi) ~ \varphi(y) ~dy d\xi$$
    $$= (2 \pi)^{-n} \int_{\R^n} \int_{\R^n} e^{i(x-y) \. \xi} ~ \xi ~ \varphi(y) ~dy d\xi$$
    $$= (2 \pi)^{-n} \int_{\R^n} e^{ix\.\xi} ~ \xi \int_{\R^n} e^{-iy \. \xi} ~ \varphi(y) ~dy d\xi$$
    $$= (2 \pi)^{-n} \int_{\R^n} e^{ix\.\xi} ~ \xi ~ \hat{\varphi}(\xi) ~d\xi$$
    $$= (2 \pi)^{-n} \int_{\R^n} D_x(e^{ix\.\xi}) ~ \hat{\varphi}(\xi) ~d\xi$$
    $$= D_x(\varphi(x)) = -i \frac{\partial}{\partial x} \varphi(x) = (\mathcal{P}\varphi)(x)$$
\end{exemple}

\begin{exemple}[Avec $\sigma (x, \xi) := x \xi$ .]
    On a
    $$(W_\sigma \varphi)(x)=(2 \pi)^{-n} \int_{\R^n} \int_{\R^n} e^{i(x-y) \. \xi}~ \sigma(\frac{x+y}{2}, ~\xi) ~ \varphi(y) ~dy d\xi$$
    $$= (2 \pi)^{-n} \int_{\R^n} \int_{\R^n} e^{i(x-y) \. \xi} ~ \frac{x+y}{2} ~\xi ~ \varphi(y) ~dy d\xi$$
    $$= \frac{1}{2}(2 \pi)^{-n} \int_{\R^n} \int_{\R^n} e^{i(x-y) \. \xi} ~ x \xi ~ \varphi(y) ~dy d\xi+\frac{1}{2}(2 \pi)^{-n} \int_{\R^n} \int_{\R^n} e^{i(x-y) \. \xi} ~ y \xi ~ \varphi(y) ~dy d\xi$$
    $$= \frac{x}{2}(2 \pi)^{-n} \int_{\R^n} ~ \xi e^{i x \. \xi}\int_{\R^n} e^{-iy \. \xi} ~ \varphi(y) ~dy d\xi+\frac{1}{2}(2 \pi)^{-n} \int_{\R^n}  \xi e^{ix \. \xi}\int_{\R^n} e^{-iy \. \xi} ~ y ~ \varphi(y) ~dy d\xi$$
    $$= \frac{x}{2}(2 \pi)^{-n} \int_{\R^n} ~ \xi e^{i x \. \xi} ~ \hat{\varphi} (\xi) d\xi-\frac{1}{2}(2 \pi)^{-n} \int_{\R^n}  \xi e^{ix \. \xi}\int_{\R^n} D_\xi(e^{-iy \. \xi})~ \varphi(y) ~dy d\xi$$
    $$= \frac{-x}{2}(2 \pi)^{-n} \int_{\R^n} ~ D_x(e^{i x \. \xi}) ~ \hat{\varphi} (\xi) d\xi-\frac{1}{2}(2 \pi)^{-n} \int_{\R^n}  \xi e^{ix \. \xi} (-x\varphi) \hat{~}(\xi) ~ d\xi$$
    $$= \frac{-x}{2} D_x(\varphi (x))-\frac{1}{2}(2 \pi)^{-n} \int_{\R^n} e^{ix \. \xi} (D_x(-x\varphi)) \hat{~}(\xi) ~ d\xi$$
    $$= \frac{-x}{2} D_x(\varphi (x))-\frac{1}{2} D_x(-x\varphi)(x)$$
    $$= \frac{ix}{2} \partial_x(\varphi (x))+\frac{i}{2} \partial_x(-x\varphi)(x) = (\frac{\hat{\mathcal{Q}}\hat{\mathcal{P}} + \hat{\mathcal{P}}\hat{\mathcal{Q}}}{2}f)(x)$$
\end{exemple}

On remarque que la transformée de Weyl a symétrisée le calcul afin de le rendre commutatif !

On aurait pu continuer en faisant les calculs pour $\sigma$ qui serait un polynôme en $\xi$ ou en $x$ ou même une fonction d'une seule variable, mais par manque de temps, je ne l'ai pas fais dans ce mémoire.\vspace{5pt}

Ces calculs prennent maintenant plus sens sous la théorie des opérateurs pseudo-différentiels.


A partir de maintenant on veut développer un calcul de la transformée de Weyl.
\subsection{Opérateurs de Hilbert-Schmidt sur $L^2(\R^n)$}
On va introduire dans ce chapitre la notion d'opérateurs de Hilbert-Schmidt sur $L^2(\R^n)$ qui sont \textbf{au coeur du développement d'un calcul associé à la transformée de Weyl}, puisque l'on verra dans la partie suivante que la notion de transformée de Weyl et d'opérateur de Hilbert-Schmidt \textbf{se confondent}.

\begin{definition}[Opérateur de Hilbert-Schmidt sur $L^2(\R^n)$]
    Soit $h \in L^2(\R^{2n})$. On définit alors \textbf{l'opérateur intégral} $S_h : L^2(\R^n) \ra L^2(\R^n)$ par
    $$(S_hf)(x) = \int_{\R^n} h(x,y)f(y)dy, ~~ x \in \R^n.$$
\end{definition}

Il est bon de vouloir munir les opérateurs de Hilbert-Schmidt sur $L^2(\R^n)$ d'opérations qui \textbf{rendent le calcul agréable}.

\begin{definition}
    Soit $f$, $g$ dans $L^2(\R^{2n})$.
    \begin{itemize}
        \item On définit $g \circ h$ sur $\R^{2n}$ par $(g \circ h)(x,y) = \int_{\R^n} g(x,z)h(z,y)dz, x, y \in \R^n$.
        \item On définit $h^*$ sur $\R^{2n}$ par $h^*(x,y) = \overline{h(y,x)}, x, y \in \R^n$.
    \end{itemize}
\end{definition}

Par la proposition suivante, on voit que \textbf{la composition est bien définie}.

\begin{proposition}
    Soit $g$ et $h$ dans $L^2(\R^{2n})$, alors $g \circ h \in L^2(\R^{2n})$.
\end{proposition}
\begin{remark}
    Cette proposition ne sera pas démontrée, la preuve utilisant essentiellement l'inégalité de Cauchy-Schwarz.
\end{remark}

On est donc amené au \textbf{théorème principal} de cette partie.

\begin{theorem}
    Soit $S_g$ et $S_h$ deux opérateurs de Hilbert-Schmidt correspondant aux noyaux $g$ et $h$ respectivement. Alors,
    \begin{itemize}
        \item $S_g S_h = S_{g \circ h}$.
        \item $S_h* = S_{h*}$.
    \end{itemize}
\end{theorem}

\begin{preuve}

    Soit $g$ et $h$ dans $L^2(\R^n)$.\vspace{5pt}

    \textbf{Preuve de $S_g S_h = S_{g \circ h}$.} \vspace{5pt}

    $$S_gS_hf(x) = S_g(S_hf)(x) = \int_{\R^n} g(x,y) S_hf(y)dy$$
    $$ = \int_{\R^n} g(x,y) \left(\int_{\R^n} h(y,z) f(z) dz\right) dy = \int_{\R^n} \int_{\R^n} g(x,y) h(y,z) f(z) dz dy$$\vspace{5pt}

    En justifiant que l'\textbf{on peut échanger les deux intégrales} par le théorème de Fubini, en justifiant que le facteur intégré est intégrable, par l'inégalité de Cauchy-Schwartz, on écrit,

    $$= \int_{\R^n} \int_{\R^n} g(x,y) h(y,z) f(z) dy dz = \int_{\R^n} \left( \int_{\R^n} g(x,y) h(y,z) dy\right) f(z) dz = \int_{\R^n} g \circ h(x,z) f(z) dz = S_{g \circ h}f(x)$$\vspace{5pt}

    \textbf{Preuve de $S_h* = S_{h*}$.}\vspace{5pt}

    $$\langle S_h f, g \rangle = \int_{\R^n} S_hf(x) \overline{g(x)} dx = \int_{\R^n} \int_{\R^n} h(x,y) f(y) dy \overline{g(x)} dx = \int_{\R^n} \int_{\R^n} h(x,y) f(y) \overline{g(x)} dy dx$$\vspace{5pt}

    Par le même argument que dans la première partie de la preuve, on échange les deux intégrales.\vspace{5pt}

    $$= \int_{\R^n} f(y) \left(\overline{\int_{\R^n}\overline{h(y,x)}g(x)}\right) dy = \int_{\R^n} f(y) \overline{S_{h*}g(y)} dy = \langle f, S_{h*}g \rangle$$\vspace{5pt}

    En conclusion $S_h$ et $S_h^*$ sont \textbf{adjoints} donc $S_h^* = S_{h^*}$.

\end{preuve}

\subsection{Produit tensoriel sur $L^2(\R^n)$}
Ce chapitre sera utile pour montrer que l'ensemble des transformées de Weyl $W_\sigma$ avec un symbole $\sigma$ dans $L^2(\R^{2n})$ est \textbf{égal} à l'ensemble des opérateurs de Hilbert-Schmidt sur $L^2(\R^n)$, c'est à dire qu'à une transformée de Weyl $W_\sigma$ on peut associer un opérateur de Hilbert-Schmidt $S_h$ et inversement. \vspace{5pt}

Au chapitre précédent, a été développé \textbf{un calcul fonctionnel sur les opérateurs de Hilbert-Schmidt}, ce qui permettra donc de munir les transformée de Weyl d'opérations naturelles.

\begin{definition}[Produit tensoriel de deux fonctions dans $L^2(\R^n)$]
    Soit $f, ~g \in L^2(\R^n)$, on définit la fonction $f \otimes g$ sur $\R^{2n}$ par $(f \otimes g)(x,y) = f(x)g(y),~~ x$, $y \in \R^n$. \vspace{5pt}

    On l'appelle le \textbf{produit tensoriel} de f et g.
\end{definition}

\begin{definition}
    Soit $f \in L^2(\R^{2n})$. On note par $\mathcal{F}_1f$ et $\mathcal{F}_2f$ les transformées de Fourier de f par rapport à la première et la seconde variable respectivement.
\end{definition}

\begin{remark}
    On note les transformées de Fourier inverses par $\mathcal{F}_1^{-1}$ et $\mathcal{F}_2^{-1}$.
\end{remark}

La proposition suivante énonce que la transformée de Fourier par rapport à une variable se comporte bien à travers le produit tensoriel.

\begin{proposition}
    Soit $f$ et $g$ dans $L^2(\R^n)$, alors,
    \begin{itemize}
        \item $\mathcal{F}_1(f \otimes g) = ( \mathcal{F}f) \otimes g$.
        \item $\mathcal{F}_2(f \otimes g) = f \otimes (\mathcal{F}g)$.
    \end{itemize}
\end{proposition}

On définit maintenant un opérateur qui sera utile pour changer de point de vue sur la façon d'appréhender la transformée de Weyl.

\begin{definition}[Opérateur de torsion]
    On définit un opérateur linéaire $T : L^2(\R^n) \ra L^2(\R^n)$ par $(T f)(x,y) = f(x+\frac{y}{2}, x - \frac{y}{2})$, $x$, $y \in \R^n$ pour tout $f$ dans $L^2(\R^{2n})$.

    On appelle $T$ \textbf{l'opérateur de torsion} sur $L^2(\R^{2n})$.
\end{definition}


On définit maintenant un opérateur linéaire $K : L^2(\R^{2n}) \ra L^2(\R^{2n})$ par
$(Kf)(x,y) = (T^{-1} \mathcal{F}_2f)(y,x)$, $x$, $y \in \R^n$ pour tout $f$ dans $L^2(\R^{2n})$.

\begin{proposition}
    L'opérateur de torsion $T : L^2(\R^n) \ra L^2(R^n)$  est \textbf{unitaire}, et
    $$(T^{-1} f)(x,y) = f\left(\frac{x+y}{2},x-y\right),~ x,~ y \in \R^n,$$
    pour tout $f \in L^2(\R^n)$.
\end{proposition}
\begin{proposition}
    L'opérateur linéaire $K$ sur $L^2(\R^n)$ a les propriétés suivantes :
    \begin{itemize}
        \item $K : L^2(\R^n) \ra L^2(\R^n)$ est un opérateur unitaire.
        \item $K = T^{-1} \mathcal{F}_2^{-1}$.
        \item $K_{\bar{f}} = (Kf)^*$, $f \in L^2(\R^{2n})$.
        \item $W(f,~g) = K^{-1}(f \otimes \bar{g})$, $f$, $g \in L^2(\R^n)$.
    \end{itemize}
\end{proposition}

\begin{remark}
    L'intérêt d'avoir introduit les deux opérateurs $T$ et $K$ sont de changer de point de vue sur les formules de la transformée de Weyl et de Wigner, afin de faciliter le calcul par la suite. \vspace{5pt}

    Le quatrième point illustre bien ce \textbf{changement de point de vue} : faire la transformée de Wigner de $f$ et $g$ revient à prendre le produit tensoriel de $f$ avec le conjugué de $g$ puis de tordre ce produit avant d'en prendre la transformée de Fourier inverse dans la seconde variable.
\end{remark}

Les théorèmes principaux de ce chapitre sont maintenant introduits.

\begin{theorem}
    Soit $\sigma \in L^2(\R^{2n})$, alors $W_\sigma : L^2(\R^n) \ra L^2(\R^n)$ est un \textbf{opérateur de Hilbert-Schmidt} de noyau $(2 \pi)^{-n/2}K_\sigma$
\end{theorem}

\begin{preuve}
    Soit $f$ et $g$ dans $L^2(\R^n)$, $\sigma$ dans $L^2(\R^{2n})$ \vspace{5pt}

    On va évaluer $W_\sigma$ par son action sur le produit scalaire.

    $$\langle W_\sigma f,~ g \rangle = (2 \pi) ^{-n/2} \int_{\R^n} \int_{\R^n}  \sigma(x, \xi) W(f,g)(x, \xi) dx d\xi$$

    On l'écrit sous forme de produit scalaire,

    $$=(2\pi)^{-n/2} \langle W(f,g),~ \bar{\sigma} \rangle$$

    D'après la proposition, on décompose $W(f,g)$,

    $$=(2 \pi)^{-n/2} \langle K^{-1}(f \otimes \bar{g}),~ \bar{\sigma} \rangle$$

    Ensuite, comme K est une \textbf{isométrie} et par la proposition, on a,
    $$= (2 \pi)^{-n/2} \langle f \otimes \bar{g},~(K_\sigma)^* \rangle = (2 \pi)^{-n/2} \int_{\R^n} \int_{\R^n} f(x) \bar{g} (y) K_\sigma (y,x) dxdy$$
    $$= (2 \pi)^{-n/2} \int_{\R^n} \left( \int_{\R^n} K_\sigma(y,x)f(x)dx \right) \bar{g}(y) dy.$$

    On reconnaît l'opérateur $S_k$ où $k(x,y) = (2 \pi)^{-n/2} K_\sigma(x,y)$, $x$, $y \in \R^n$ définit sur $\R^{2n}$.
    $$= \int_{\R^n}(S_kf)(y) \bar{g}(y)dy = \langle S_kf,~g\rangle.$$
\end{preuve}

\begin{theorem}
    Soit $S_h$ un opérateur de Hilbert-Schmidt dans $L^2(\R^n)$, alors $S_h$ est une \textbf{transformée de Weyl} de symbole $\sigma = (2 \pi)^{n/2}K^{-1}h$.
\end{theorem}

\begin{preuve}
    Soit $S_h$ un opérateur de Hilbert-Schmidt, $h \in L^2(\R^{2n})$.\vspace{5pt}

    En posant $\sigma = (2 \pi)^{n/2} K^{-1}h$, d'après le théorème précédent, $W_\sigma$ est un opérateur de Hilbert-Schmidt de noyau $(2\pi)^{-n/2} K (2 \pi)^{n/2} K^{-1} h = h$
\end{preuve}
En conclusion l'espace des opérateurs de Hilbert-Schmidt sur $L^2(\R^n)$ \textbf{est égal} à l'espace des transformées de Weyl avec symbole dans $L^2(\R^n)$.
\subsection{$\mathcal{H}^*$-algèbre et calcul de Weyl}

Ce chapitre final va développer \textbf{un calcul fonctionnel pour les transformées de Weyl} avec un symbole dans $L^2(\R^n)$. \vspace{5pt}

Maintenant que l'on a identifié les transformées de Weyl avec symbole dans $L^2(\R^n)$ aux opérateurs de Hilbert-Schmidt dans $L^2(\R^n)$. On va essayer de développer la notion de $\mathcal{H}^*$-algèbre. Au coeur de cette partie.

\begin{definition}
    Soit $\mathcal{H}$ un espace de Hilbert \textbf{séparable complexe}, avec $\langle \.,~ \. \rangle$ son produit scalaire et $\| \. \|$ sa norme. \vspace{5pt}

    On munit $\mathcal{H}$ de deux opérations :
    $a,~ b \mapsto ab$ et $a \mapsto a^*$.\vspace{5pt}

    On dit que $\mathcal{H}$ est une \textbf{$\mathcal{H^*}$}-algèbre si les deux opérations satisfont :


    \begin{itemize}\setlength{\itemsep}{1pt}
        \item $ a(b+c) = ab+ac$, $(a+b)c = ac+bc$.
        \item $ \lambda (ab) = (\lambda a) b = a (\lambda b)$.
        \item $ a(bc) = (ab)c$.
        \item $a^{**}=a$.
        \item $(a+b)^* = a^*+b^*$.
        \item $(ab)^* = b^* a^*$.
        \item $ (\lambda a)^* = \bar{\lambda} a^*$.
        \item $\| a^* \| = \| a \|$.
        \item $\|ab\| \leq \|a\| \|b\|$.
        \item $\langle ab, c \rangle = \langle b, a^*c \rangle$.
    \end{itemize}

\end{definition}

On cherche à partir de maintenant de munir les ensembles que l'on connaît de nouvelles structures, différentes et plus riche, afin de faire un lien profond entre ces objets.

\begin{proposition}
    L'espace de Hilbert $L^2(\R^{2n})$ équipé des deux opérations $f,~ g \mapsto f \circ g$ et $f \mapsto f^*$ avec l'opération $\circ$ définit dans le chapitre sur les opérateurs de Hilbert-Schmidt est une $\mathcal{H}^*$-algèbre. \vspace{5pt}

    On note la $\mathcal{H}^*$-algèbre obtenue par $(L^2(\R^{2n}), \circ, *)$.
\end{proposition}


\begin{definition}% l'introduire avant

    On note $HS(L^2(\R^n))$ \textbf{l'ensemble des opérateurs de Hilbert-Schmidt} sur $L^2(\R^n)$.\vspace{5pt}

    On le munit d'un produit scalaire définit pour tout $A$, $B \in HS(L^2(\R^n))$ par $(A,B)_{HS} = \overset{\infty}{\underset{k=1}{\sum}} \langle A_{\varphi_k, B_{\varphi_k}} \rangle$,\vspace{5pt}

    où les $\varphi_k$ forment une base orthonormale de $L^2(\R^n)$.
\end{definition}

\begin{proposition}
    L'espace $HS(L^2(\R^n))$ équipé avec le produit scalaire $(~,~)_{HS}$ et les opérations de produit usuel et d'adjoint est une $\mathcal{H}^*$-algèbre.\vspace{5pt}

    On la note encore $HS(L^2(\R^n))$.

\end{proposition}

\begin{proposition}
    La $\mathcal{H}^*$-algèbre$(L^2(\R^{2n}), \circ, *)$ est isométriquement $*-$isomorphe à la $\mathcal{H}^*$-algèbre $HS(L^2(\R^n))$ par l'application $h \mapsto S_h$.
\end{proposition}
\begin{remark}
    C'est à dire qu'il existe un isomorphisme isométrique.
\end{remark}

\begin{definition}
    Soit $f$, $g \in L^2(\R^n)$, on définit la fonction $f \times g$ sur $\R^{2n}$ par
    $$f \times g = K^{-1} (Kf \circ Kg)$$
\end{definition}

\begin{proposition}
    L'espace de Hilbert $L^2(\R^{2n})$ équipé des deux opérations $f,~ g \mapsto f \times g$ et $f \mapsto \bar{f}$ est une $\mathcal{H}^*$-algèbre.\vspace{5pt}

    De plus $K$ réalise un \textbf{isomorphisme} entre $(L^2(\R^{2n}), \times, -)$ et $(L^2(\R^{2n}), \circ, *)$.
\end{proposition}

\begin{theorem}
    L'application
    $$L^2(\R^{2n}) \ni \sigma \mapsto (2 \pi)^{n/2} W_\sigma \in HS(L^2(\R^n))$$ est un $*-$isomorphisme isométrique de $(L^2(\R^{2n}), \times, -)$ dans $HS(L^2(\R^{2n}))$. \vspace{5pt}

    En conséquence, on a,
    \begin{itemize}
        \item $W_\sigma^* = W_{\bar{\sigma}}$
        \item $W_\sigma W_\tau = W_{(2 \pi)^{-n/2} \sigma \times \tau}$
        \item $W_\sigma + W_\tau = W_{\sigma+\tau}$
        \item $\lambda W_\sigma = W_{\lambda \sigma}$
        \item $\|W_\sigma\|_* \leq (2 \pi)^{-n/2} \|\sigma\|_{L^2(\R^n)} = \|W_\sigma\|_{HS}$
    \end{itemize}
\end{theorem}

\begin{preuve}
    On a le diagramme suivant, qui montre facilement le théorème.

    \begin{center}
        \begin{tikzpicture}[x=3cm,y=2cm, >=latex, ->]
            \node(1) at(0:-2){$(L^2(\R^{2n}), \times, -)$};
            \node(2) at(0:0){$(L^2(\R^{2n}), \circ, *)$};
            \node(3) at(-90:1){$HS(L^2(\R^{2n}))$};

            \foreach \k/\l in {1/3} {\draw (\k)--(\l) node[midway, above, sloped]{};}
            \foreach \k/\l in {2/3} {\draw (\k)--(\l) node[midway, right]{S};}
            \foreach \k/\l in {1/2} {\draw (\k)--(\l) node[midway, above, sloped]{K};}
        \end{tikzpicture}
    \end{center}
\end{preuve}

L'une des conséquences les plus importante de ce théorème est que $W_\sigma$ est auto-adjoint si et seulement si $\sigma$ est à valeurs réelles. \vspace{10pt}

En conclusion, on a crée un cadre formel et bien défini pour la transformée de Weyl, qu'on a construite pas à pas, en motivant les définitions introduites, puis on a fini par en donner quelques exemples. On a ensuite fini le mémoire par donner une structure opérationnelle fonctionnelle sur les transformées de Weyl en faisant la correspondance avec les opérateurs de Hilbert-Schmidt. J'aurais aimé réussir à faire un lien plus profond entre la transformée de Wigner et la transformée de Weyl ainsi que faire plus de calculs d'illustrations. Ce mémoire m'a permis d'avoir plus de rigueur mathématique dans ma formulation des choses, grâce à l'aide de mon encadrant, j'ai pu comprendre que je n'étais pas assez rigoureux ou précis. J'ai principalement utilisé les sources \cite{weyl} et \cite{folland}

\printbibliography

\end{document}
